%=====================================================================
%  We Know Group | Kato Kiosk Guide  (AM + IT)
%  Command Centre setup, silent app install, kiosk lockdown,
%  and (Appendix A) the FreedomPay FCC machine build.
%  Shared styling lives in wkg-preamble.tex.
%=====================================================================
\documentclass[a4paper,10pt]{article}
%=====================================================================
%  We Know Group | Shared preamble for the Kato kiosk guides
%
%  Branding per "We Know Group Brand Guidelines V2 2023":
%    - Primary colours: Graphite #121212, Azure #00A6D3, White
%    - Secondary: Sky Blue #E5F1F9, Sea Blue #008BB2, Night Blue #0C2752
%    - Greys: Night #575A5C, Mid #ADADAD, Light #F4F4F4
%    - Typeface: Circular Pro (approximated here with Inter, the licensed
%      brand stand-in used in the WeKnow Employee Handbook)
%    - Logos are the official WeKnow assets in images/ -- never redrawn.
%
%  Both Kato kiosk guides share this file so their branding stays in step.
%
%  Usage, from a document that has already declared \documentclass:
%      %=====================================================================
%  We Know Group | Shared preamble for the Kato kiosk guides
%
%  Branding per "We Know Group Brand Guidelines V2 2023":
%    - Primary colours: Graphite #121212, Azure #00A6D3, White
%    - Secondary: Sky Blue #E5F1F9, Sea Blue #008BB2, Night Blue #0C2752
%    - Greys: Night #575A5C, Mid #ADADAD, Light #F4F4F4
%    - Typeface: Circular Pro (approximated here with Inter, the licensed
%      brand stand-in used in the WeKnow Employee Handbook)
%    - Logos are the official WeKnow assets in images/ -- never redrawn.
%
%  Both Kato kiosk guides share this file so their branding stays in step.
%
%  Usage, from a document that has already declared \documentclass:
%      %=====================================================================
%  We Know Group | Shared preamble for the Kato kiosk guides
%
%  Branding per "We Know Group Brand Guidelines V2 2023":
%    - Primary colours: Graphite #121212, Azure #00A6D3, White
%    - Secondary: Sky Blue #E5F1F9, Sea Blue #008BB2, Night Blue #0C2752
%    - Greys: Night #575A5C, Mid #ADADAD, Light #F4F4F4
%    - Typeface: Circular Pro (approximated here with Inter, the licensed
%      brand stand-in used in the WeKnow Employee Handbook)
%    - Logos are the official WeKnow assets in images/ -- never redrawn.
%
%  Both Kato kiosk guides share this file so their branding stays in step.
%
%  Usage, from a document that has already declared \documentclass:
%      \input{wkg-preamble}
%      \renewcommand{\guidetitle}{Short title for the page footer}
%=====================================================================

%---------------------------------------------------------------------
% Fonts & encoding -- Inter (brand stand-in for Circular Pro)
%---------------------------------------------------------------------
\usepackage[T1]{fontenc}
\usepackage[utf8]{inputenc}
\usepackage[sfdefault]{inter}
\renewcommand{\familydefault}{\sfdefault}
\usepackage{microtype}

%---------------------------------------------------------------------
% Layout
%---------------------------------------------------------------------
\usepackage[a4paper,margin=20mm,top=28mm,bottom=22mm,headheight=20pt,headsep=14pt]{geometry}
\usepackage{graphicx}
\graphicspath{{images/}}
\usepackage[table]{xcolor}
\usepackage{tikz}
\usetikzlibrary{arrows.meta,positioning,calc,shapes.misc,fit,backgrounds}
\usepackage{tcolorbox}
\tcbuselibrary{skins,breakable}
\usepackage{booktabs}
\usepackage{array}
\usepackage{tabularx}
\usepackage{enumitem}
\usepackage{fancyhdr}
\usepackage{titlesec}
\usepackage{listings}
\usepackage{float}
\usepackage{lastpage}
\usepackage{amssymb}
\usepackage{caption}
\usepackage[hidelinks,colorlinks=true,urlcolor=weknowsea,linkcolor=weknownight]{hyperref}

%---------------------------------------------------------------------
% We Know Group palette (Brand Guidelines V2 2023, Section 3)
%---------------------------------------------------------------------
\definecolor{weknowazure}{HTML}{00A6D3}     % Hero colour
\definecolor{weknowgraphite}{HTML}{121212}
\definecolor{weknowsky}{HTML}{E5F1F9}
\definecolor{weknowsea}{HTML}{008BB2}
\definecolor{weknownight}{HTML}{0C2752}
\definecolor{weknowgreyN}{HTML}{575A5C}
\definecolor{weknowgreyM}{HTML}{ADADAD}
\definecolor{weknowgreyL}{HTML}{F4F4F4}
\definecolor{weknowazure80}{HTML}{33B8DC}    % Azure tints (brand: tint Azure only)
\definecolor{weknowazure60}{HTML}{66CAE5}
\definecolor{weknowazure40}{HTML}{99DBED}
\definecolor{weknowazure20}{HTML}{CCEDF6}

%---------------------------------------------------------------------
% Headers / footers -- \guidetitle is set by each document
%---------------------------------------------------------------------
\providecommand{\guidetitle}{Kato Kiosk Guide}
\pagestyle{fancy}
\fancyhf{}
% \sectionmark uses \markright, so the running head must read \rightmark.
\fancyhead[L]{\footnotesize\color{weknowgreyN}\nouppercase{\rightmark}}
\fancyhead[R]{\raisebox{-0.35\height}{\includegraphics[height=15pt]{wkg-logo-horizontal}}}
\renewcommand{\headrulewidth}{0.8pt}
\renewcommand{\headrule}{\hbox to\headwidth{\color{weknowazure}\leaders\hrule height \headrulewidth\hfill}}
\fancyfoot[L]{\scriptsize\color{weknowgreyN}Internal --- We Know Group}
\fancyfoot[C]{\scriptsize\color{weknowgreyN}\guidetitle}
\fancyfoot[R]{\scriptsize\color{weknowgreyN}Page \thepage\ of \pageref{LastPage}}
\renewcommand{\footrulewidth}{0.4pt}
\renewcommand{\footrule}{\hbox to\headwidth{\color{weknowgreyM}\leaders\hrule height \footrulewidth\hfill}}

%---------------------------------------------------------------------
% Section styling
%---------------------------------------------------------------------
\titleformat{\section}
  {\color{weknowgraphite}\LARGE\bfseries}
  {\color{weknowazure}\thesection}{0.7em}{}
  [{\vspace{3pt}\color{weknowazure}\titlerule[1.4pt]}]
\titleformat{\subsection}
  {\color{weknowsea}\Large\bfseries}
  {\color{weknowazure}\thesubsection}{0.6em}{}
\titlespacing*{\section}{0pt}{20pt}{12pt}
\titlespacing*{\subsection}{0pt}{16pt}{6pt}
% The running head takes its text from \wksecmark, not from the section
% title, so role badges never leak into the header. \wksection and
% \wkrolesection (defined below) set it. Subsections must not overwrite
% the section in the head.
\newcommand{\wksecmark}{}
\renewcommand{\sectionmark}[1]{\markright{\thesection\quad\wksecmark}}
\renewcommand{\subsectionmark}[1]{}

%---------------------------------------------------------------------
% Captions
%---------------------------------------------------------------------
\captionsetup{font=small,labelfont={bf,color=weknownight},
              textfont={color=weknowgreyN},skip=5pt}

%---------------------------------------------------------------------
% Screenshot helper -- \screenshot[width fraction]{file}{caption}
%---------------------------------------------------------------------
\newcommand{\screenshot}[3][0.72]{%
  \begin{figure}[H]
    \centering
    \setlength{\fboxsep}{0pt}\setlength{\fboxrule}{0.6pt}%
    {\color{weknowgreyM}\fbox{\includegraphics[width=#1\linewidth]{#2}}}%
    \caption{#3}
  \end{figure}}

%---------------------------------------------------------------------
% Callout boxes
%---------------------------------------------------------------------
\newtcolorbox{wkinfo}[1][]{%
  enhanced,breakable,colback=weknowsky,colframe=weknowazure,
  boxrule=0pt,leftrule=2.5pt,arc=1.5pt,
  fonttitle=\bfseries,coltitle=weknownight,title={#1},
  left=6pt,right=6pt,top=5pt,bottom=5pt}
\newtcolorbox{wkwarn}[1][]{%
  enhanced,breakable,colback=weknowgreyL,colframe=weknownight,
  boxrule=0pt,leftrule=2.5pt,arc=1.5pt,
  fonttitle=\bfseries,coltitle=white,title={#1},
  left=6pt,right=6pt,top=5pt,bottom=5pt}

%---------------------------------------------------------------------
% Code listings
%---------------------------------------------------------------------
\usepackage{upquote}      % straight ' and ` in listings, so copied code runs
\usepackage{attachfile2}  % embed scripts in the PDF, see \textattachfile
\lstdefinestyle{wkcode}{
  upquote=true,
  basicstyle=\ttfamily\footnotesize\color{weknowgraphite},
  backgroundcolor=\color{weknowgreyL},
  frame=single, framerule=0pt, rulecolor=\color{weknowgreyL},
  breaklines=true, breakatwhitespace=false,
  columns=fullflexible, keepspaces=true,
  aboveskip=8pt, belowskip=8pt, xleftmargin=6pt, framexleftmargin=6pt}
\lstset{style=wkcode}

\newcolumntype{L}[1]{>{\raggedright\arraybackslash}p{#1}}
\setlist{nosep,itemsep=4pt}

%---------------------------------------------------------------------
% Page-breaking behaviour
%   - never strand a single line at the top or bottom of a page
%   - \wkneed{N\baselineskip} makes a page break attractive when less
%     than N lines remain, so a heading cannot land at the page foot
%     with nothing under it (needspace's trick, inlined to avoid a
%     dependency)
%---------------------------------------------------------------------
\widowpenalty=10000
\clubpenalty=10000
\displaywidowpenalty=10000
\makeatletter
\newcommand{\wkneed}[1]{%
  \begingroup
    \setlength{\dimen@}{#1}%
    \vskip\z@\@plus\dimen@
    \penalty-100
    \vskip\z@\@plus-\dimen@
  \endgroup}
\makeatother

%---------------------------------------------------------------------
% Shorthands
%---------------------------------------------------------------------
\newcommand{\wpath}[1]{{\ttfamily\small #1}}
\newcommand{\tick}{{\color{weknowsea}$\checkmark$}}
\newcommand{\emptybox}{{\color{weknowgreyM}$\square$}}

% Role badges -- white on solid Azure, in two sizes.
%   \role     for the section heading (set \LARGE, so the badge sits at
%             \normalsize and reads as a tag, not a second title)
%   \roletoc  for the table of contents (set \normalsize, so the badge
%             drops to \scriptsize)
% Do not call these directly in a heading -- use \wksection or
% \wkrolesection below, which pick the right size per context.
% Robust: these are written to the .toc file, so they must not expand there.
% #1 size, #2 padding, #3 baseline raise, #4 text
\DeclareRobustCommand{\wkbadge}[4]{%
  {\setlength{\fboxsep}{#2}%
   \raisebox{#3}{%
     \colorbox{weknowazure}{\textcolor{white}{#1\bfseries #4}}}}}
\DeclareRobustCommand{\role}[1]{\texorpdfstring{%
  \hspace{0.45em}\wkbadge{\normalsize}{3.5pt}{-0.35ex}{#1}}{[#1]}}
\DeclareRobustCommand{\roletoc}[1]{\texorpdfstring{%
  \hspace{0.35em}\wkbadge{\scriptsize}{2pt}{-0.25ex}{#1}}{[#1]}}

%---------------------------------------------------------------------
% Section wrappers -- these set the running head and, where a role is
% given, place the badge at heading size in the body and at the smaller
% size in the table of contents.
%   \wksection{Final checklist}
%   \wkrolesection{Part 2 --- Install and lock the kiosk}{IT}
%---------------------------------------------------------------------
\newcommand{\wksection}[1]{%
  \renewcommand{\wksecmark}{#1}%
  \wkneed{6\baselineskip}%
  \section{#1}}
\newcommand{\wkrolesection}[2]{%
  \renewcommand{\wksecmark}{#1}%
  \wkneed{6\baselineskip}%
  \section[#1 \roletoc{#2}]{#1 \role{#2}}}

%---------------------------------------------------------------------
% Cover page -- \makewkgcover{title}{subtitle}{audience}
%---------------------------------------------------------------------
\newcommand{\makewkgcover}[3]{%
\begin{titlepage}
\thispagestyle{empty}
\begin{tikzpicture}[remember picture,overlay]
  % --- 'Wink' echo: symmetric Azure-tint rounded squares framing the top ---
  \fill[weknowazure20,rounded corners=8mm]
       ([xshift=-46mm,yshift=-46mm]current page.north east) rectangle
       ([xshift=14mm,yshift=-8mm]current page.north east);
  \fill[weknowazure40,rounded corners=5mm]
       ([xshift=-30mm,yshift=-40mm]current page.north east) rectangle
       ([xshift=8mm,yshift=-13mm]current page.north east);
  \fill[weknowazure20,rounded corners=8mm]
       ([xshift=-14mm,yshift=-46mm]current page.north west) rectangle
       ([xshift=46mm,yshift=-8mm]current page.north west);
  \fill[weknowazure40,rounded corners=5mm]
       ([xshift=-8mm,yshift=-40mm]current page.north west) rectangle
       ([xshift=30mm,yshift=-13mm]current page.north west);

  % --- Large centred WeKnow logo (stacked variant) = the hero ---
  \node[anchor=north] at ([yshift=-50mm]current page.north)
       {\includegraphics[width=52mm]{wkg-logo-stacked}};

  % --- Centred title / Azure hero rule / subtitle ---
  \node[anchor=north,align=center,text width=185mm]
       at ([yshift=-126mm]current page.north)
       {{\color{weknowgraphite}\fontsize{26}{31}\selectfont\bfseries #1}\\[11pt]
        {\color{weknowazure}\rule{42mm}{2pt}}\\[11pt]
        {\color{weknowsea}\fontsize{15}{19}\selectfont #2}};

  % --- Metadata panel ---
  \node[anchor=north, fill=weknowsky, rounded corners=3mm,
        inner xsep=12mm, inner ysep=8mm]
       at ([yshift=-178mm]current page.north)
       {\renewcommand{\arraystretch}{1.5}%
        \begin{tabular}{@{}l@{\hspace{12mm}}l@{}}
          {\color{weknowsea}\bfseries Version} & {\color{weknowgraphite}1.0}\\
          {\color{weknowsea}\bfseries Date} & {\color{weknowgraphite}August 2026}\\
          {\color{weknowsea}\bfseries Classification} & {\color{weknowgraphite}Internal}\\
          {\color{weknowsea}\bfseries Prepared by} & {\color{weknowgraphite}Vedant Kalbag}\\
          {\color{weknowsea}\bfseries For} & {\color{weknowgraphite}#3}\\
        \end{tabular}};

  % --- Confidentiality notice ---
  \node[anchor=south,align=center,text width=150mm]
       at ([yshift=20mm]current page.south)
       {{\color{weknowgreyM}\rule{55mm}{0.4pt}}\\[8pt]
        {\color{weknowgreyN}\footnotesize This document contains confidential and proprietary
         information belonging to We Know Group. It is intended solely for internal use and may
         not be reproduced or distributed without prior written consent.}};
\end{tikzpicture}
\end{titlepage}}

%      \renewcommand{\guidetitle}{Short title for the page footer}
%=====================================================================

%---------------------------------------------------------------------
% Fonts & encoding -- Inter (brand stand-in for Circular Pro)
%---------------------------------------------------------------------
\usepackage[T1]{fontenc}
\usepackage[utf8]{inputenc}
\usepackage[sfdefault]{inter}
\renewcommand{\familydefault}{\sfdefault}
\usepackage{microtype}

%---------------------------------------------------------------------
% Layout
%---------------------------------------------------------------------
\usepackage[a4paper,margin=20mm,top=28mm,bottom=22mm,headheight=20pt,headsep=14pt]{geometry}
\usepackage{graphicx}
\graphicspath{{images/}}
\usepackage[table]{xcolor}
\usepackage{tikz}
\usetikzlibrary{arrows.meta,positioning,calc,shapes.misc,fit,backgrounds}
\usepackage{tcolorbox}
\tcbuselibrary{skins,breakable}
\usepackage{booktabs}
\usepackage{array}
\usepackage{tabularx}
\usepackage{enumitem}
\usepackage{fancyhdr}
\usepackage{titlesec}
\usepackage{listings}
\usepackage{float}
\usepackage{lastpage}
\usepackage{amssymb}
\usepackage{caption}
\usepackage[hidelinks,colorlinks=true,urlcolor=weknowsea,linkcolor=weknownight]{hyperref}

%---------------------------------------------------------------------
% We Know Group palette (Brand Guidelines V2 2023, Section 3)
%---------------------------------------------------------------------
\definecolor{weknowazure}{HTML}{00A6D3}     % Hero colour
\definecolor{weknowgraphite}{HTML}{121212}
\definecolor{weknowsky}{HTML}{E5F1F9}
\definecolor{weknowsea}{HTML}{008BB2}
\definecolor{weknownight}{HTML}{0C2752}
\definecolor{weknowgreyN}{HTML}{575A5C}
\definecolor{weknowgreyM}{HTML}{ADADAD}
\definecolor{weknowgreyL}{HTML}{F4F4F4}
\definecolor{weknowazure80}{HTML}{33B8DC}    % Azure tints (brand: tint Azure only)
\definecolor{weknowazure60}{HTML}{66CAE5}
\definecolor{weknowazure40}{HTML}{99DBED}
\definecolor{weknowazure20}{HTML}{CCEDF6}

%---------------------------------------------------------------------
% Headers / footers -- \guidetitle is set by each document
%---------------------------------------------------------------------
\providecommand{\guidetitle}{Kato Kiosk Guide}
\pagestyle{fancy}
\fancyhf{}
% \sectionmark uses \markright, so the running head must read \rightmark.
\fancyhead[L]{\footnotesize\color{weknowgreyN}\nouppercase{\rightmark}}
\fancyhead[R]{\raisebox{-0.35\height}{\includegraphics[height=15pt]{wkg-logo-horizontal}}}
\renewcommand{\headrulewidth}{0.8pt}
\renewcommand{\headrule}{\hbox to\headwidth{\color{weknowazure}\leaders\hrule height \headrulewidth\hfill}}
\fancyfoot[L]{\scriptsize\color{weknowgreyN}Internal --- We Know Group}
\fancyfoot[C]{\scriptsize\color{weknowgreyN}\guidetitle}
\fancyfoot[R]{\scriptsize\color{weknowgreyN}Page \thepage\ of \pageref{LastPage}}
\renewcommand{\footrulewidth}{0.4pt}
\renewcommand{\footrule}{\hbox to\headwidth{\color{weknowgreyM}\leaders\hrule height \footrulewidth\hfill}}

%---------------------------------------------------------------------
% Section styling
%---------------------------------------------------------------------
\titleformat{\section}
  {\color{weknowgraphite}\LARGE\bfseries}
  {\color{weknowazure}\thesection}{0.7em}{}
  [{\vspace{3pt}\color{weknowazure}\titlerule[1.4pt]}]
\titleformat{\subsection}
  {\color{weknowsea}\Large\bfseries}
  {\color{weknowazure}\thesubsection}{0.6em}{}
\titlespacing*{\section}{0pt}{20pt}{12pt}
\titlespacing*{\subsection}{0pt}{16pt}{6pt}
% The running head takes its text from \wksecmark, not from the section
% title, so role badges never leak into the header. \wksection and
% \wkrolesection (defined below) set it. Subsections must not overwrite
% the section in the head.
\newcommand{\wksecmark}{}
\renewcommand{\sectionmark}[1]{\markright{\thesection\quad\wksecmark}}
\renewcommand{\subsectionmark}[1]{}

%---------------------------------------------------------------------
% Captions
%---------------------------------------------------------------------
\captionsetup{font=small,labelfont={bf,color=weknownight},
              textfont={color=weknowgreyN},skip=5pt}

%---------------------------------------------------------------------
% Screenshot helper -- \screenshot[width fraction]{file}{caption}
%---------------------------------------------------------------------
\newcommand{\screenshot}[3][0.72]{%
  \begin{figure}[H]
    \centering
    \setlength{\fboxsep}{0pt}\setlength{\fboxrule}{0.6pt}%
    {\color{weknowgreyM}\fbox{\includegraphics[width=#1\linewidth]{#2}}}%
    \caption{#3}
  \end{figure}}

%---------------------------------------------------------------------
% Callout boxes
%---------------------------------------------------------------------
\newtcolorbox{wkinfo}[1][]{%
  enhanced,breakable,colback=weknowsky,colframe=weknowazure,
  boxrule=0pt,leftrule=2.5pt,arc=1.5pt,
  fonttitle=\bfseries,coltitle=weknownight,title={#1},
  left=6pt,right=6pt,top=5pt,bottom=5pt}
\newtcolorbox{wkwarn}[1][]{%
  enhanced,breakable,colback=weknowgreyL,colframe=weknownight,
  boxrule=0pt,leftrule=2.5pt,arc=1.5pt,
  fonttitle=\bfseries,coltitle=white,title={#1},
  left=6pt,right=6pt,top=5pt,bottom=5pt}

%---------------------------------------------------------------------
% Code listings
%---------------------------------------------------------------------
\usepackage{upquote}      % straight ' and ` in listings, so copied code runs
\usepackage{attachfile2}  % embed scripts in the PDF, see \textattachfile
\lstdefinestyle{wkcode}{
  upquote=true,
  basicstyle=\ttfamily\footnotesize\color{weknowgraphite},
  backgroundcolor=\color{weknowgreyL},
  frame=single, framerule=0pt, rulecolor=\color{weknowgreyL},
  breaklines=true, breakatwhitespace=false,
  columns=fullflexible, keepspaces=true,
  aboveskip=8pt, belowskip=8pt, xleftmargin=6pt, framexleftmargin=6pt}
\lstset{style=wkcode}

\newcolumntype{L}[1]{>{\raggedright\arraybackslash}p{#1}}
\setlist{nosep,itemsep=4pt}

%---------------------------------------------------------------------
% Page-breaking behaviour
%   - never strand a single line at the top or bottom of a page
%   - \wkneed{N\baselineskip} makes a page break attractive when less
%     than N lines remain, so a heading cannot land at the page foot
%     with nothing under it (needspace's trick, inlined to avoid a
%     dependency)
%---------------------------------------------------------------------
\widowpenalty=10000
\clubpenalty=10000
\displaywidowpenalty=10000
\makeatletter
\newcommand{\wkneed}[1]{%
  \begingroup
    \setlength{\dimen@}{#1}%
    \vskip\z@\@plus\dimen@
    \penalty-100
    \vskip\z@\@plus-\dimen@
  \endgroup}
\makeatother

%---------------------------------------------------------------------
% Shorthands
%---------------------------------------------------------------------
\newcommand{\wpath}[1]{{\ttfamily\small #1}}
\newcommand{\tick}{{\color{weknowsea}$\checkmark$}}
\newcommand{\emptybox}{{\color{weknowgreyM}$\square$}}

% Role badges -- white on solid Azure, in two sizes.
%   \role     for the section heading (set \LARGE, so the badge sits at
%             \normalsize and reads as a tag, not a second title)
%   \roletoc  for the table of contents (set \normalsize, so the badge
%             drops to \scriptsize)
% Do not call these directly in a heading -- use \wksection or
% \wkrolesection below, which pick the right size per context.
% Robust: these are written to the .toc file, so they must not expand there.
% #1 size, #2 padding, #3 baseline raise, #4 text
\DeclareRobustCommand{\wkbadge}[4]{%
  {\setlength{\fboxsep}{#2}%
   \raisebox{#3}{%
     \colorbox{weknowazure}{\textcolor{white}{#1\bfseries #4}}}}}
\DeclareRobustCommand{\role}[1]{\texorpdfstring{%
  \hspace{0.45em}\wkbadge{\normalsize}{3.5pt}{-0.35ex}{#1}}{[#1]}}
\DeclareRobustCommand{\roletoc}[1]{\texorpdfstring{%
  \hspace{0.35em}\wkbadge{\scriptsize}{2pt}{-0.25ex}{#1}}{[#1]}}

%---------------------------------------------------------------------
% Section wrappers -- these set the running head and, where a role is
% given, place the badge at heading size in the body and at the smaller
% size in the table of contents.
%   \wksection{Final checklist}
%   \wkrolesection{Part 2 --- Install and lock the kiosk}{IT}
%---------------------------------------------------------------------
\newcommand{\wksection}[1]{%
  \renewcommand{\wksecmark}{#1}%
  \wkneed{6\baselineskip}%
  \section{#1}}
\newcommand{\wkrolesection}[2]{%
  \renewcommand{\wksecmark}{#1}%
  \wkneed{6\baselineskip}%
  \section[#1 \roletoc{#2}]{#1 \role{#2}}}

%---------------------------------------------------------------------
% Cover page -- \makewkgcover{title}{subtitle}{audience}
%---------------------------------------------------------------------
\newcommand{\makewkgcover}[3]{%
\begin{titlepage}
\thispagestyle{empty}
\begin{tikzpicture}[remember picture,overlay]
  % --- 'Wink' echo: symmetric Azure-tint rounded squares framing the top ---
  \fill[weknowazure20,rounded corners=8mm]
       ([xshift=-46mm,yshift=-46mm]current page.north east) rectangle
       ([xshift=14mm,yshift=-8mm]current page.north east);
  \fill[weknowazure40,rounded corners=5mm]
       ([xshift=-30mm,yshift=-40mm]current page.north east) rectangle
       ([xshift=8mm,yshift=-13mm]current page.north east);
  \fill[weknowazure20,rounded corners=8mm]
       ([xshift=-14mm,yshift=-46mm]current page.north west) rectangle
       ([xshift=46mm,yshift=-8mm]current page.north west);
  \fill[weknowazure40,rounded corners=5mm]
       ([xshift=-8mm,yshift=-40mm]current page.north west) rectangle
       ([xshift=30mm,yshift=-13mm]current page.north west);

  % --- Large centred WeKnow logo (stacked variant) = the hero ---
  \node[anchor=north] at ([yshift=-50mm]current page.north)
       {\includegraphics[width=52mm]{wkg-logo-stacked}};

  % --- Centred title / Azure hero rule / subtitle ---
  \node[anchor=north,align=center,text width=185mm]
       at ([yshift=-126mm]current page.north)
       {{\color{weknowgraphite}\fontsize{26}{31}\selectfont\bfseries #1}\\[11pt]
        {\color{weknowazure}\rule{42mm}{2pt}}\\[11pt]
        {\color{weknowsea}\fontsize{15}{19}\selectfont #2}};

  % --- Metadata panel ---
  \node[anchor=north, fill=weknowsky, rounded corners=3mm,
        inner xsep=12mm, inner ysep=8mm]
       at ([yshift=-178mm]current page.north)
       {\renewcommand{\arraystretch}{1.5}%
        \begin{tabular}{@{}l@{\hspace{12mm}}l@{}}
          {\color{weknowsea}\bfseries Version} & {\color{weknowgraphite}1.0}\\
          {\color{weknowsea}\bfseries Date} & {\color{weknowgraphite}August 2026}\\
          {\color{weknowsea}\bfseries Classification} & {\color{weknowgraphite}Internal}\\
          {\color{weknowsea}\bfseries Prepared by} & {\color{weknowgraphite}Vedant Kalbag}\\
          {\color{weknowsea}\bfseries For} & {\color{weknowgraphite}#3}\\
        \end{tabular}};

  % --- Confidentiality notice ---
  \node[anchor=south,align=center,text width=150mm]
       at ([yshift=20mm]current page.south)
       {{\color{weknowgreyM}\rule{55mm}{0.4pt}}\\[8pt]
        {\color{weknowgreyN}\footnotesize This document contains confidential and proprietary
         information belonging to We Know Group. It is intended solely for internal use and may
         not be reproduced or distributed without prior written consent.}};
\end{tikzpicture}
\end{titlepage}}

%      \renewcommand{\guidetitle}{Short title for the page footer}
%=====================================================================

%---------------------------------------------------------------------
% Fonts & encoding -- Inter (brand stand-in for Circular Pro)
%---------------------------------------------------------------------
\usepackage[T1]{fontenc}
\usepackage[utf8]{inputenc}
\usepackage[sfdefault]{inter}
\renewcommand{\familydefault}{\sfdefault}
\usepackage{microtype}

%---------------------------------------------------------------------
% Layout
%---------------------------------------------------------------------
\usepackage[a4paper,margin=20mm,top=28mm,bottom=22mm,headheight=20pt,headsep=14pt]{geometry}
\usepackage{graphicx}
\graphicspath{{images/}}
\usepackage[table]{xcolor}
\usepackage{tikz}
\usetikzlibrary{arrows.meta,positioning,calc,shapes.misc,fit,backgrounds}
\usepackage{tcolorbox}
\tcbuselibrary{skins,breakable}
\usepackage{booktabs}
\usepackage{array}
\usepackage{tabularx}
\usepackage{enumitem}
\usepackage{fancyhdr}
\usepackage{titlesec}
\usepackage{listings}
\usepackage{float}
\usepackage{lastpage}
\usepackage{amssymb}
\usepackage{caption}
\usepackage[hidelinks,colorlinks=true,urlcolor=weknowsea,linkcolor=weknownight]{hyperref}

%---------------------------------------------------------------------
% We Know Group palette (Brand Guidelines V2 2023, Section 3)
%---------------------------------------------------------------------
\definecolor{weknowazure}{HTML}{00A6D3}     % Hero colour
\definecolor{weknowgraphite}{HTML}{121212}
\definecolor{weknowsky}{HTML}{E5F1F9}
\definecolor{weknowsea}{HTML}{008BB2}
\definecolor{weknownight}{HTML}{0C2752}
\definecolor{weknowgreyN}{HTML}{575A5C}
\definecolor{weknowgreyM}{HTML}{ADADAD}
\definecolor{weknowgreyL}{HTML}{F4F4F4}
\definecolor{weknowazure80}{HTML}{33B8DC}    % Azure tints (brand: tint Azure only)
\definecolor{weknowazure60}{HTML}{66CAE5}
\definecolor{weknowazure40}{HTML}{99DBED}
\definecolor{weknowazure20}{HTML}{CCEDF6}

%---------------------------------------------------------------------
% Headers / footers -- \guidetitle is set by each document
%---------------------------------------------------------------------
\providecommand{\guidetitle}{Kato Kiosk Guide}
\pagestyle{fancy}
\fancyhf{}
% \sectionmark uses \markright, so the running head must read \rightmark.
\fancyhead[L]{\footnotesize\color{weknowgreyN}\nouppercase{\rightmark}}
\fancyhead[R]{\raisebox{-0.35\height}{\includegraphics[height=15pt]{wkg-logo-horizontal}}}
\renewcommand{\headrulewidth}{0.8pt}
\renewcommand{\headrule}{\hbox to\headwidth{\color{weknowazure}\leaders\hrule height \headrulewidth\hfill}}
\fancyfoot[L]{\scriptsize\color{weknowgreyN}Internal --- We Know Group}
\fancyfoot[C]{\scriptsize\color{weknowgreyN}\guidetitle}
\fancyfoot[R]{\scriptsize\color{weknowgreyN}Page \thepage\ of \pageref{LastPage}}
\renewcommand{\footrulewidth}{0.4pt}
\renewcommand{\footrule}{\hbox to\headwidth{\color{weknowgreyM}\leaders\hrule height \footrulewidth\hfill}}

%---------------------------------------------------------------------
% Section styling
%---------------------------------------------------------------------
\titleformat{\section}
  {\color{weknowgraphite}\LARGE\bfseries}
  {\color{weknowazure}\thesection}{0.7em}{}
  [{\vspace{3pt}\color{weknowazure}\titlerule[1.4pt]}]
\titleformat{\subsection}
  {\color{weknowsea}\Large\bfseries}
  {\color{weknowazure}\thesubsection}{0.6em}{}
\titlespacing*{\section}{0pt}{20pt}{12pt}
\titlespacing*{\subsection}{0pt}{16pt}{6pt}
% The running head takes its text from \wksecmark, not from the section
% title, so role badges never leak into the header. \wksection and
% \wkrolesection (defined below) set it. Subsections must not overwrite
% the section in the head.
\newcommand{\wksecmark}{}
\renewcommand{\sectionmark}[1]{\markright{\thesection\quad\wksecmark}}
\renewcommand{\subsectionmark}[1]{}

%---------------------------------------------------------------------
% Captions
%---------------------------------------------------------------------
\captionsetup{font=small,labelfont={bf,color=weknownight},
              textfont={color=weknowgreyN},skip=5pt}

%---------------------------------------------------------------------
% Screenshot helper -- \screenshot[width fraction]{file}{caption}
%---------------------------------------------------------------------
\newcommand{\screenshot}[3][0.72]{%
  \begin{figure}[H]
    \centering
    \setlength{\fboxsep}{0pt}\setlength{\fboxrule}{0.6pt}%
    {\color{weknowgreyM}\fbox{\includegraphics[width=#1\linewidth]{#2}}}%
    \caption{#3}
  \end{figure}}

%---------------------------------------------------------------------
% Callout boxes
%---------------------------------------------------------------------
\newtcolorbox{wkinfo}[1][]{%
  enhanced,breakable,colback=weknowsky,colframe=weknowazure,
  boxrule=0pt,leftrule=2.5pt,arc=1.5pt,
  fonttitle=\bfseries,coltitle=weknownight,title={#1},
  left=6pt,right=6pt,top=5pt,bottom=5pt}
\newtcolorbox{wkwarn}[1][]{%
  enhanced,breakable,colback=weknowgreyL,colframe=weknownight,
  boxrule=0pt,leftrule=2.5pt,arc=1.5pt,
  fonttitle=\bfseries,coltitle=white,title={#1},
  left=6pt,right=6pt,top=5pt,bottom=5pt}

%---------------------------------------------------------------------
% Code listings
%---------------------------------------------------------------------
\usepackage{upquote}      % straight ' and ` in listings, so copied code runs
\usepackage{attachfile2}  % embed scripts in the PDF, see \textattachfile
\lstdefinestyle{wkcode}{
  upquote=true,
  basicstyle=\ttfamily\footnotesize\color{weknowgraphite},
  backgroundcolor=\color{weknowgreyL},
  frame=single, framerule=0pt, rulecolor=\color{weknowgreyL},
  breaklines=true, breakatwhitespace=false,
  columns=fullflexible, keepspaces=true,
  aboveskip=8pt, belowskip=8pt, xleftmargin=6pt, framexleftmargin=6pt}
\lstset{style=wkcode}

\newcolumntype{L}[1]{>{\raggedright\arraybackslash}p{#1}}
\setlist{nosep,itemsep=4pt}

%---------------------------------------------------------------------
% Page-breaking behaviour
%   - never strand a single line at the top or bottom of a page
%   - \wkneed{N\baselineskip} makes a page break attractive when less
%     than N lines remain, so a heading cannot land at the page foot
%     with nothing under it (needspace's trick, inlined to avoid a
%     dependency)
%---------------------------------------------------------------------
\widowpenalty=10000
\clubpenalty=10000
\displaywidowpenalty=10000
\makeatletter
\newcommand{\wkneed}[1]{%
  \begingroup
    \setlength{\dimen@}{#1}%
    \vskip\z@\@plus\dimen@
    \penalty-100
    \vskip\z@\@plus-\dimen@
  \endgroup}
\makeatother

%---------------------------------------------------------------------
% Shorthands
%---------------------------------------------------------------------
\newcommand{\wpath}[1]{{\ttfamily\small #1}}
\newcommand{\tick}{{\color{weknowsea}$\checkmark$}}
\newcommand{\emptybox}{{\color{weknowgreyM}$\square$}}

% Role badges -- white on solid Azure, in two sizes.
%   \role     for the section heading (set \LARGE, so the badge sits at
%             \normalsize and reads as a tag, not a second title)
%   \roletoc  for the table of contents (set \normalsize, so the badge
%             drops to \scriptsize)
% Do not call these directly in a heading -- use \wksection or
% \wkrolesection below, which pick the right size per context.
% Robust: these are written to the .toc file, so they must not expand there.
% #1 size, #2 padding, #3 baseline raise, #4 text
\DeclareRobustCommand{\wkbadge}[4]{%
  {\setlength{\fboxsep}{#2}%
   \raisebox{#3}{%
     \colorbox{weknowazure}{\textcolor{white}{#1\bfseries #4}}}}}
\DeclareRobustCommand{\role}[1]{\texorpdfstring{%
  \hspace{0.45em}\wkbadge{\normalsize}{3.5pt}{-0.35ex}{#1}}{[#1]}}
\DeclareRobustCommand{\roletoc}[1]{\texorpdfstring{%
  \hspace{0.35em}\wkbadge{\scriptsize}{2pt}{-0.25ex}{#1}}{[#1]}}

%---------------------------------------------------------------------
% Section wrappers -- these set the running head and, where a role is
% given, place the badge at heading size in the body and at the smaller
% size in the table of contents.
%   \wksection{Final checklist}
%   \wkrolesection{Part 2 --- Install and lock the kiosk}{IT}
%---------------------------------------------------------------------
\newcommand{\wksection}[1]{%
  \renewcommand{\wksecmark}{#1}%
  \wkneed{6\baselineskip}%
  \section{#1}}
\newcommand{\wkrolesection}[2]{%
  \renewcommand{\wksecmark}{#1}%
  \wkneed{6\baselineskip}%
  \section[#1 \roletoc{#2}]{#1 \role{#2}}}

%---------------------------------------------------------------------
% Cover page -- \makewkgcover{title}{subtitle}{audience}
%---------------------------------------------------------------------
\newcommand{\makewkgcover}[3]{%
\begin{titlepage}
\thispagestyle{empty}
\begin{tikzpicture}[remember picture,overlay]
  % --- 'Wink' echo: symmetric Azure-tint rounded squares framing the top ---
  \fill[weknowazure20,rounded corners=8mm]
       ([xshift=-46mm,yshift=-46mm]current page.north east) rectangle
       ([xshift=14mm,yshift=-8mm]current page.north east);
  \fill[weknowazure40,rounded corners=5mm]
       ([xshift=-30mm,yshift=-40mm]current page.north east) rectangle
       ([xshift=8mm,yshift=-13mm]current page.north east);
  \fill[weknowazure20,rounded corners=8mm]
       ([xshift=-14mm,yshift=-46mm]current page.north west) rectangle
       ([xshift=46mm,yshift=-8mm]current page.north west);
  \fill[weknowazure40,rounded corners=5mm]
       ([xshift=-8mm,yshift=-40mm]current page.north west) rectangle
       ([xshift=30mm,yshift=-13mm]current page.north west);

  % --- Large centred WeKnow logo (stacked variant) = the hero ---
  \node[anchor=north] at ([yshift=-50mm]current page.north)
       {\includegraphics[width=52mm]{wkg-logo-stacked}};

  % --- Centred title / Azure hero rule / subtitle ---
  \node[anchor=north,align=center,text width=185mm]
       at ([yshift=-126mm]current page.north)
       {{\color{weknowgraphite}\fontsize{26}{31}\selectfont\bfseries #1}\\[11pt]
        {\color{weknowazure}\rule{42mm}{2pt}}\\[11pt]
        {\color{weknowsea}\fontsize{15}{19}\selectfont #2}};

  % --- Metadata panel ---
  \node[anchor=north, fill=weknowsky, rounded corners=3mm,
        inner xsep=12mm, inner ysep=8mm]
       at ([yshift=-178mm]current page.north)
       {\renewcommand{\arraystretch}{1.5}%
        \begin{tabular}{@{}l@{\hspace{12mm}}l@{}}
          {\color{weknowsea}\bfseries Version} & {\color{weknowgraphite}1.0}\\
          {\color{weknowsea}\bfseries Date} & {\color{weknowgraphite}August 2026}\\
          {\color{weknowsea}\bfseries Classification} & {\color{weknowgraphite}Internal}\\
          {\color{weknowsea}\bfseries Prepared by} & {\color{weknowgraphite}Vedant Kalbag}\\
          {\color{weknowsea}\bfseries For} & {\color{weknowgraphite}#3}\\
        \end{tabular}};

  % --- Confidentiality notice ---
  \node[anchor=south,align=center,text width=150mm]
       at ([yshift=20mm]current page.south)
       {{\color{weknowgreyM}\rule{55mm}{0.4pt}}\\[8pt]
        {\color{weknowgreyN}\footnotesize This document contains confidential and proprietary
         information belonging to We Know Group. It is intended solely for internal use and may
         not be reproduced or distributed without prior written consent.}};
\end{tikzpicture}
\end{titlepage}}

\renewcommand{\guidetitle}{Kato Kiosk Guide --- AM \& IT}

\begin{document}

\makewkgcover{Kato: \ \ Kiosk \ Installation\\ \& \ Configuration \ Guide}
             {Command Centre, Machine Build, \\App Install \& Lockdown}
             {Account Managers \& IT}

%=====================================================================
% HOW TO USE
%=====================================================================
\section*{How to use this guide}
\markright{How to use this guide}

This guide covers \textbf{installing and configuring a kiosk for a specific
installation}, from building the machine to locking it down. It is split
by who does the work:

\begin{center}
\renewcommand{\arraystretch}{1.35}\small
\begin{tabularx}{\textwidth}{L{22mm} L{30mm} X}
\rowcolor{weknownight}\color{white}\textbf{Section} &
\color{white}\textbf{Who} & \color{white}\textbf{What}\\
\S\ref{sec:before} & Both &
Before you start: the location information to collect, the two
identifiers, and the important paths on the machine.\\
\rowcolor{weknowgreyL}
\S\ref{sec:hierarchy} & Account Manager &
How configuration resolves: which level of the hierarchy each setting
belongs on.\\
\S\ref{sec:part1} & Account Manager &
Part 1: create the location, kiosk and installation record in Command
Centre, clear the prep checks, and send IT the kiosk's link.\\
\rowcolor{weknowgreyL}
\S\ref{sec:part2} & IT &
Part 2: copy the Kiosk ID and generate the API key, then one N-central
setup script builds the machine, installs the Kato app and locks it into
kiosk mode. Or do it by hand.\\
\S\ref{sec:checklist} & Both &
Final checklist, run before leaving site.\\
\rowcolor{weknowgreyL}
\S\ref{sec:trouble} & IT &
When things go wrong: common faults, the setup script's messages, the
identity file checks, and finding the printer's COM port.\\
Appendix~\ref{sec:fcc} & IT &
Building the machine by hand: install and verify FreedomPay FCC. Once per
machine. The setup script does this as its first step.\\
\rowcolor{weknowgreyL}
Appendix~\ref{sec:quickref} & Both &
One-page quick reference for someone who has done this before.\\
\end{tabularx}
\end{center}

\begin{figure}[H]
\centering
\begin{tikzpicture}[
  box/.style={draw=weknowazure,fill=white,rounded corners=3pt,
              text width=26mm,minimum height=12mm,align=center,
              inner sep=3pt,font=\scriptsize},
  lane/.style={font=\small\bfseries,text=weknowsea,align=center,
               text width=20mm},
  arr/.style={-{Stealth[length=2.4mm]},thick,weknowazure},
  handoff/.style={-{Stealth[length=2.4mm]},thick,weknowgreyN,dashed}]

  % --- Account Manager lane ---
  \node[box]  (t1) at (0,0)   {Create location\\and kiosk};
  \node[box]  (t2) at (3.5,0) {Configure and\\clear prep checks};
  \node[box]  (t3) at (7.0,0) {Create installation\\record};
  \node[box]  (t4) at (10.5,0) {Send IT the\\kiosk link};
  \node[lane] (l1) at (-2.8,0) {Part 1\\Account Manager};

  % --- IT lane ---
  \node[box]  (b1) at (3.5,-3.0) {Copy Kiosk ID,\\generate API key};
  \node[box]  (b2) at (7.0,-3.0) {Install the\\Kato app};
  \node[box]  (b3) at (10.5,-3.0) {Lock into\\kiosk mode};
  \node[lane] (l2) at (-2.8,-3.0) {Part 2\\IT};

  % --- Lane bands ---
  \begin{scope}[on background layer]
    \node[fill=weknowsky,rounded corners=4pt,inner sep=3.5mm,
          fit=(l1)(t1)(t4)] {};
    \node[fill=weknowgreyL,rounded corners=4pt,inner sep=3.5mm,
          fit=(l2)(b1)(b3)] {};
  \end{scope}

  \draw[arr] (t1) -- (t2);
  \draw[arr] (t2) -- (t3);
  \draw[arr] (t3) -- (t4);
  \draw[arr] (b1) -- (b2);
  \draw[arr] (b2) -- (b3);

  \draw[handoff] (t4.south) -- ++(0,-0.75) -| (b1.north);
  \node[font=\scriptsize,text=weknowgreyN,anchor=south]
       at ([yshift=-0.75cm]$(t4.south)!0.5!(t4.south -| b1.north)$) {Kiosk link};
\end{tikzpicture}
\caption{The Account Manager configures the kiosk and sends IT its Command Centre link. IT take the Kiosk ID and a fresh API key from there, install the app and lock the machine down.}
\end{figure}

\begin{wkwarn}[The machine must be built first]
FreedomPay FCC and the PAL packages must be installed and verified on the
machine before the kiosk goes live. IT do this once per machine: the Part~2
setup script does it as its first step, or by hand with
Appendix~\ref{sec:fcc}. The prep checks in Part~1 cannot honestly be
cleared until it is done.
\end{wkwarn}

\begin{wkinfo}[Golden rules]
\begin{itemize}
  \item The \textbf{API key is shown once}. IT generate it at the moment
        of install and paste it straight into the setup script or install
        command --- it is never saved or sent anywhere. If it is lost, generate a new one.
  \item The \textbf{Kiosk ID} and the \textbf{Workstation ID} are
        different things --- see \S\ref{sec:ids}.
  \item Do not clear a prep check you have not actually verified. It is
        the only thing standing between a half-built kiosk and a live
        one.
\end{itemize}
\end{wkinfo}

\clearpage
\tableofcontents
\clearpage

%=====================================================================
% 1  BEFORE YOU START
%=====================================================================
\wksection{Before you start}\label{sec:before}

\subsection{Minimum information needed for a location}

The Account Manager must have all of the following for the location
before an installation can be created. If any are missing, chase them
first.

\begin{center}
\renewcommand{\arraystretch}{1.3}\small
\begin{tabularx}{\textwidth}{X X}
\rowcolor{weknowgreyL}
\emptybox\ Hotel name & \emptybox\ Currency\\
\emptybox\ Which products to enable & \emptybox\ Timezone\\
\rowcolor{weknowgreyL}
\emptybox\ Products for the home screen & \emptybox\ Relevant airports\\
\emptybox\ Default language (kiosk + PDQ) & \emptybox\ Default Uber suggestions\\
\rowcolor{weknowgreyL}
\emptybox\ VAT / Tax ID & \emptybox\ Tax entity name\\
\emptybox\ Banking details (in full) & \\
\end{tabularx}
\end{center}

This is the \emph{minimum}. Many other parameters can be configured at
each level of the hierarchy --- see \S\ref{sec:hierarchy} for where each
one belongs.

\begin{wkwarn}[Legacy fields --- not required]
\textbf{RPS customer code} and \textbf{outlet code} are legacy fields.
Do not chase the site for them.
\end{wkwarn}

\subsection{Two identifiers, two different things}\label{sec:ids}

\begin{center}
\renewcommand{\arraystretch}{1.35}\small
\begin{tabularx}{\textwidth}{L{34mm} X}
\rowcolor{weknownight}\color{white}\textbf{Identifier} &
\color{white}\textbf{What it is}\\
\textbf{Kiosk ID} &
A UUID generated by Command Centre when the kiosk record is created. The
app uses it to identify itself to the server.\\
\rowcolor{weknowgreyL}
\textbf{Workstation ID} &
The physical machine label, e.g.\ \wpath{WKG-POS-617}. Also entered as the
kiosk's serial number in Command Centre.\\
\end{tabularx}
\end{center}

Both go into the install command in Part~2 --- do not swap them.

\subsection{Important paths}\label{sec:paths}

Every file and setting this guide touches on the kiosk machine.

\begin{center}
\renewcommand{\arraystretch}{1.3}\small
\begin{tabularx}{\textwidth}{L{34mm} X}
\rowcolor{weknownight}\color{white}\textbf{What} &
\color{white}\textbf{Where}\\
Identity file &
\wpath{C:\textbackslash Users\textbackslash SelfBooker\textbackslash AppData\textbackslash Roaming\textbackslash WKG Kiosk\textbackslash kiosk-identity.json}\\
\rowcolor{weknowgreyL}
Kato app &
\wpath{C:\textbackslash Users\textbackslash SelfBooker\textbackslash AppData\textbackslash Local\textbackslash Programs\textbackslash Kiosk\textbackslash Kiosk.exe}\\
Installer drop folder &
\wpath{C:\textbackslash Users\textbackslash SelfBooker\textbackslash Documents}\\
\rowcolor{weknowgreyL}
Shell Launcher utility &
\wpath{C:\textbackslash Kiosk\textbackslash SSM25\_files\textbackslash Shell-Launcher-Utilities\textbackslash SLUtility.exe}\\
FCC logs &
\wpath{C:\textbackslash ProgramData\textbackslash FreedomPay\textbackslash Freeway Commerce Connect\textbackslash log}\\
\rowcolor{weknowgreyL}
FCC registry key &
\wpath{HKEY\_LOCAL\_MACHINE\textbackslash SOFTWARE\textbackslash Freedompay\textbackslash Atlas\textbackslash PAL}\\
Latest app release &
\href{https://github.com/WeKnowTechnology/wkg-electron-app/releases}{\wpath{github.com/WeKnowTechnology/wkg-electron-app/releases}}\\
\rowcolor{weknowgreyL}
Setup scripts &
\href{https://github.com/vedant-kalbag-wkg/kato-installation-guide/tree/main/scripts/v2}{\wpath{github.com/vedant-kalbag-wkg/kato-installation-guide/tree/main/scripts/v2}}\\
Setup script log &
\wpath{C:\textbackslash installation.log}\\
\rowcolor{weknowgreyL}
Step 1's PAL record &
\wpath{HKEY\_LOCAL\_MACHINE\textbackslash SOFTWARE\textbackslash WeKnow\textbackslash Kiosk}, value \wpath{PalPackage}\\
\end{tabularx}
\end{center}

%=====================================================================
% 2  CONFIGURATION HIERARCHY
%=====================================================================
\wkrolesection{How configuration resolves}{Account Manager}\label{sec:hierarchy}

Configuration is layered. An individual kiosk should carry the
\textbf{minimum configuration needed} --- only the things genuinely unique
to that kiosk. Everything else resolves down from the levels above it.

\begin{figure}[H]
\centering
\begin{tikzpicture}[
  lvl/.style={rounded corners=2.5pt,draw=weknowazure60,
              minimum height=8mm,align=center,inner xsep=4pt,
              font=\scriptsize,text=weknownight},
  node distance=1.6mm]

  \node[lvl,fill=weknowgreyL]    (g)  {Global};
  \node[lvl,fill=weknowazure20,right=of g]  (co) {Country};
  \node[lvl,fill=weknowazure20,right=of co] (re) {Region\\(State)};
  \node[lvl,fill=weknowazure40,right=of re] (ci) {City};
  \node[lvl,fill=weknowazure40,right=of ci] (lg) {Location\\Group};
  \node[lvl,fill=weknowazure60,right=of lg] (lo) {Location};
  \node[lvl,fill=weknowazure60,right=of lo] (kg) {Kiosk\\Group};
  \node[lvl,fill=weknowazure80,right=of kg] (ki) {Kiosk};

  \draw[-{Stealth[length=2.4mm]},thick,weknowazure]
       ([yshift=-5mm]g.south west) -- ([yshift=-5mm]ki.south east);
  \node[font=\scriptsize,text=weknowgreyN,anchor=north]
       at ([yshift=-6mm]$(g.south west)!0.5!(ki.south east)$)
       {widest scope \quad$\longrightarrow$\quad most specific --- and highest priority};
\end{tikzpicture}
\caption{The configuration hierarchy. Anything set further right overrides the same parameter set further left.}
\end{figure}

\subsection{The rule}

A parameter set at a \textbf{more granular} level overrides the same
parameter set at a \textbf{wider} level. So a currency set on a Country
applies to every location in it, until a specific Location --- or a single
Kiosk --- sets its own.

\begin{wkinfo}[Set each parameter at the widest level where it holds true]
If every kiosk in a country shares a currency, set it once on the Country;
do not repeat it on each kiosk. Repeating a value further down does no
harm on day one, but it is the thing that makes later changes miss ---
you update the Country and the kiosk keeps its own stale override.
\end{wkinfo}

\subsection{When a kiosk is in more than one group}\label{sec:groupsorder}

A kiosk can belong to \textbf{several kiosk groups at once}. Assign them
on the kiosk's \textbf{Edit} tab, under \textbf{Group assignments}. A
location can likewise belong to several location groups, and both tiers
follow the same rule below.

Assigned groups carry a \textbf{display order}, and that order decides
conflicts: where two of a kiosk's groups set the same parameter, the group
\textbf{lower in the list wins} and overrides the ones above it. The list
runs widest first, exactly like the hierarchy above it --- later beats
earlier --- so the bottom of the list is the most specific position, not
the top.

\begin{wkinfo}[Order the groups deliberately]
Put the most specific group at the bottom. A kiosk that is both
\emph{eSIM-only} and part of a wider \emph{Heathrow} group should carry
the narrower group below the broader one, so the narrow group's settings
are the ones that survive.

A newly assigned group is added at the bottom, which makes it the winner
by default. If you assign a broad group after a narrow one, move it up.
\end{wkinfo}

\begin{wkwarn}[The screen does not tell you which end wins]
The rows show only an order number. Nothing on the panel says whether the
top or the bottom takes precedence, so after any reorder confirm the
outcome on the kiosk's \textbf{Resolved} tab rather than trusting the
list.
\end{wkwarn}

\subsection{Working out where a parameter belongs}

Ask how widely the value is true, and set it at that level:

\begin{center}
\renewcommand{\arraystretch}{1.35}\small
\begin{tabularx}{\textwidth}{L{34mm} X}
\rowcolor{weknownight}\color{white}\textbf{Level} &
\color{white}\textbf{Typically carries}\\
Global &
Anything true of every WeKnow kiosk everywhere. Rarely touched.\\
\rowcolor{weknowgreyL}
Country &
Things that follow the market --- currency, tax treatment, default
language.\\
Region (State) &
Regional variations within a country, where they exist.\\
\rowcolor{weknowgreyL}
City &
Things that follow geography --- relevant airports, default Uber
suggestions.\\
Location Group &
A set of sites treated the same, e.g.\ one hotel brand or one operator
across several properties.\\
\rowcolor{weknowgreyL}
Location &
The property itself --- hotel name, address, tax entity, banking details.\\
Kiosk Group &
A set of kiosks that share a specific configuration. \textbf{Not limited
to one location} --- a kiosk group can span sites, e.g.\ every Uber-only,
HEx-only or eSIM-only kiosk across the estate. Also used for kiosks
treated the same within a single site.\\
\rowcolor{weknowgreyL}
Kiosk &
Only what is unique to the single machine.\\
\end{tabularx}
\end{center}

The right-hand column is a guide to intent, not a fixed schema --- many
parameters can legitimately be set at more than one level. What matters is
choosing the widest one that is still correct.

\begin{wkinfo}[Check what a kiosk actually ends up with]
The kiosk's \textbf{Resolved} tab shows the effective configuration after
the whole hierarchy has been applied. Use it to confirm a kiosk is picking
up what you intended --- and to find the stale override when it is not.
\end{wkinfo}

%=====================================================================
% 3  PART 1 -- COMMAND CENTRE
%=====================================================================
\wkrolesection{Part 1 --- Command Centre}{Account Manager}\label{sec:part1}

Done from any machine with a browser. You are creating and configuring the
kiosk record, clearing its prep checks, scheduling the installation, and
sending IT a link to the kiosk. You do \textbf{not} generate the API key ---
IT do that at install time.

\subsection{Create the location and the kiosk}

\begin{enumerate}
  \item Open \href{https://command.weknow.technology}{Command Centre}
        and sign in. Ask the WKG Technology team if you do not have an
        account.

  \item Go to \href{https://command.weknow.technology/locations}{Locations}
        and create the location, working down the hierarchy:
        \textbf{Country $\rightarrow$ Region $\rightarrow$ City
        $\rightarrow$ Location}. Open it and \textbf{copy the Location
        ID} --- you need it when you create the installation record.
\end{enumerate}

\screenshot[0.50]{cc-location-id}{The location Overview tab. Use the copy icon next to \textbf{Location ID}.}

\begin{enumerate}\setcounter{enumi}{2}
  \item \emph{Only if required:} create a group from
        \href{https://command.weknow.technology/groups}{Groups}. Groups
        apply shared configuration to several sites or several kiosks at
        once:
    \begin{itemize}
      \item \textbf{Location groups} --- sites treated the same, e.g.\ one
            hotel brand or one operator across several properties.
      \item \textbf{Kiosk groups} --- kiosks that share a specific
            configuration. These are \emph{not} limited to a single
            location: a kiosk group can gather every Uber-only, HEx-only
            or eSIM-only kiosk across different sites.
    \end{itemize}
    Skip this for a single standalone kiosk.

  \item Go to \textbf{Kiosks} and select \textbf{Add Kiosk}:
    \begin{itemize}
      \item \textbf{Name} --- how the kiosk is referred to operationally,
            e.g.\ \emph{T2/T3 Tunnel - 2}.
      \item \textbf{Serial number} --- the workstation label,
            e.g.\ \wpath{WKG-POS-617}.
      \item \textbf{Hardware model} --- e.g.\ \emph{K10/20}.
      \item \textbf{Firmware version} --- e.g.\ \emph{1.0.0}.
    \end{itemize}
    Then select \textbf{Create Kiosk}.
\end{enumerate}

\screenshot[0.42]{cc-add-kiosk}{Add Kiosk. The kiosk is created unassigned --- it is tied to a location by the installation record.}

\begin{enumerate}\setcounter{enumi}{4}
  \item The new kiosk appears as \textbf{In Stock}. That is expected ---
        it stays In Stock until it is installed.

  \item If the kiosk belongs to any kiosk groups, open its \textbf{Edit}
        tab and assign them under \textbf{Group assignments}. A kiosk can
        be in several groups at once, and the order they sit in decides
        which one wins a conflict --- see \S\ref{sec:groupsorder}.
\end{enumerate}

\screenshot[0.36]{cc-group-assignments}{Group assignments on the kiosk Edit tab. Pick a group and select \textbf{Assign}; assigned groups are held in a display order, widest first, and the group at the bottom wins a conflict.}

\subsection{Set the payment credentials}

Set the FreedomPay \textbf{Store ID (StID)} and \textbf{Terminal ID
(TID)} in the kiosk configuration, so payment requests are signed for the
correct merchant.

\subsection{Clear the prep checks}

Open the kiosk and go to the \textbf{Edit} tab. Mark the prep checks
\textbf{Complete} to confirm the kiosk is ready to config and install.

\screenshot[0.74]{cc-prep-checks}{Prep checks on the kiosk Edit tab.}

Before clearing them, verify with IT that:

\begin{itemize}
  \item all the hardware is plugged in and working --- printer, PDQ,
        sensors and scanners;
  \item FreedomPay has been installed, and the
        \wpath{FreewayClientService} and \wpath{FreewayServerService} are
        running (step~1 of the setup script, or Appendix~\ref{sec:fcc}).
\end{itemize}

\subsection{Create the installation record}

The installation record ties the kiosk to its location and moves it out of
\emph{In Stock}.

\begin{enumerate}
  \item Go to
        \href{https://command.weknow.technology/installations}{Installations}
        and select \textbf{Add Installation}.

  \item Complete the form:
    \begin{itemize}
      \item \textbf{Kiosk} --- the kiosk you created above.
      \item \textbf{Location ID} --- the UUID copied earlier.
      \item \textbf{Planned date} --- the installation date.
      \item \textbf{Hardware model} and \textbf{Serial number} --- must
            match the kiosk record.
      \item \textbf{Installer name}, \textbf{Installer contact} and
            \textbf{Installation company}.
      \item \textbf{Notes} --- anything the engineer on site needs to
            know, e.g.\ \emph{``Install near main entrance''}.
    \end{itemize}

  \item Select \textbf{Create Installation}.
\end{enumerate}

\screenshot[0.36]{cc-add-installation}{Add Installation --- links the kiosk to its location and records who is installing it.}

\subsection{Send IT the kiosk link}\label{sec:handover}

Open the kiosk in Command Centre and copy the page's address from the
browser bar. Send it to IT in the install ticket or chat --- the link is
\textbf{not sensitive} and needs no special handling. Everything IT need
is on that page: the Kiosk ID, the Workstation ID (the kiosk's serial
number), and the API keys tab.

If the printer is on a port other than the standard \wpath{COM8}, say so
in the same message.

\begin{wkinfo}[Done when]
The kiosk record exists and is configured, its prep checks are cleared,
the installation record is created, and IT have the link to the kiosk.
\end{wkinfo}

%=====================================================================
% 3  PART 2 -- IT
%=====================================================================
\wkrolesection{Part 2 --- Install and lock the kiosk}{IT}\label{sec:part2}

You need the link to the kiosk in Command Centre from the Account Manager,
and a Command Centre account --- ask the WKG Technology team if you do not
have one.

There are two ways to do this part. Use \textbf{Option A}. Use Option~B
only if the script cannot be used, or to finish by hand a job the script
stopped on.

\begin{center}
\renewcommand{\arraystretch}{1.35}\small
\begin{tabularx}{\textwidth}{L{44mm} X}
\rowcolor{weknownight}\color{white}\textbf{Option} &
\color{white}\textbf{What it does}\\
A --- Setup script from N-central (\S\ref{sec:ncentral}) &
One paste into the N-central remote connection. It builds the machine
(Appendix~\ref{sec:fcc}), installs the Kato app, locks the machine into
kiosk mode and restarts it into Kato, checking each result as it goes.\\
\rowcolor{weknowgreyL}
B --- By hand (\S\ref{sec:manual}) &
Appendix~\ref{sec:fcc}, then the silent install signed in as
\wpath{SelfBooker}, then the lockdown with \wpath{SLUtility}.\\
\end{tabularx}
\end{center}

\subsection{Before you install}\label{sec:preinstall}

\begin{itemize}
  \item The \textbf{PDQ is plugged in}. The build cannot be verified
        without it.
  \item You know the kiosk's \textbf{region}: \wpath{UK}, \wpath{DE},
        \wpath{ES} or \wpath{PT}. It picks the FreedomPay config and PAL
        package.
  \item \wpath{SelfBooker} is \textbf{signed in}, or \textbf{signs in
        automatically} when the machine starts. The script installs the
        app as \wpath{SelfBooker} either way.
  \item For Option~A: the machine's \textbf{computer name is its
        Workstation ID}, e.g.\ \wpath{WKG-POS-617} --- the kiosk's serial
        number in Command Centre. The script uses the computer name as the
        Workstation ID. If they differ, use Option~B, where you type the
        Workstation ID yourself.
\end{itemize}

\subsection{Get the Kiosk ID and API key}\label{sec:getkeys}

Do this \textbf{immediately before you install}, with the setup script or
install command open and ready to paste into.

\begin{enumerate}
  \item Open the kiosk link from the Account Manager. On the
        \textbf{Overview} tab, \textbf{copy the Kiosk ID}. The kiosk's
        \textbf{serial number} is the Workstation ID.
\end{enumerate}

\screenshot[0.48]{cc-kiosk-id}{Kiosk Overview --- copy the \textbf{Kiosk ID}.}

\begin{enumerate}\setcounter{enumi}{1}
  \item Switch to the \textbf{API keys} tab and select \textbf{Generate
        API key}. This is the key the kiosk uses to authenticate to the
        server.
\end{enumerate}

\screenshot[0.46]{cc-api-key}{Kiosk API keys. Generate the key here.}

\begin{wkwarn}[Use the API key straight away --- do not save it]
The full key is displayed \textbf{only once, at creation}. Paste it
directly into the setup script or install command. Do not save it anywhere
or send it to anyone --- not the ticket, chat, email or a notes file. If
you lose it before the install is done, simply generate a new one.
\end{wkwarn}

\subsection{Option A --- Setup script from N-central}\label{sec:ncentral}

The setup script, \wpath{setup\_kiosk.ps1}, runs in the N-central remote
connection's PowerShell, as the machine's system account. It holds three
steps and runs them in order, stopping at the first one that fails.

\begin{enumerate}
  \item Open the N-central remote connection to the machine. In the
        \textbf{System Shell} panel, select the \textbf{PowerShell} icon at
        the top right to switch the shell to PowerShell mode.
\end{enumerate}

\screenshot[0.94]{ncentral-powershell}{N-central System Shell. The PowerShell icon is the first of the three at the top right; the shell shows \textbf{Powershell} once switched.}

\begin{enumerate}\setcounter{enumi}{1}
  \item On your own computer, open the script with the button below. It
        opens as text in the browser. Select all (\textbf{Ctrl+A}), copy
        (\textbf{Ctrl+C}) and paste it into a new \textbf{Notepad} window.
        \begin{center}
        \href{https://raw.githubusercontent.com/vedant-kalbag-wkg/kato-installation-guide/main/scripts/v2/setup_kiosk.ps1}{%
          \colorbox{weknowazure}{\color{white}\bfseries\strut\quad Open setup\_kiosk.ps1\quad}}
        \end{center}
  \item In Notepad, replace the three values at the top ---
        \wpath{\$Region}, \wpath{\$KioskId} and \wpath{\$ApiKey} --- with
        the region and the values from \S\ref{sec:getkeys}. Keep the
        quotes, and leave everything below them alone. The script stops
        straight away if a value is still a placeholder.
  \item Select all and copy again, then close Notepad \textbf{without
        saving} --- it holds the API key.
  \item Paste the whole script into the N-central shell. It runs as it
        arrives. If the last line is left waiting at the prompt, press
        \textbf{Enter}.
  \item Watch the output. Each step ends with \texttt{=== SUCCESS ===} and
        a summary, and the run ends with
        \texttt{=== ALL 3 STEPS PASSED ===}. The machine then restarts and
        comes up in the Kato app as \wpath{SelfBooker}, with no desktop.
\end{enumerate}

\begin{center}
\renewcommand{\arraystretch}{1.35}\small
\begin{tabularx}{\textwidth}{L{20mm} X L{40mm}}
\rowcolor{weknownight}\color{white}\textbf{Step} &
\color{white}\textbf{What it does} &
\color{white}\textbf{Skipped on a rerun}\\
1 FreedomPay &
Appendix~\ref{sec:fcc} for the region: installs FCC 5.3.2.15 silently,
deletes the \wpath{Atlas\textbackslash PAL} key, copies the config, runs
the PAL package and restarts the services in order. Then waits up to two
minutes for \emph{Lane Open}, \emph{Successfully pushed} and
\emph{PalWriteSucceeded} in the FCC log. &
The FCC install, if FCC is there. The PAL package and the service
restart, if this region's package was the last one run.\\
\rowcolor{weknowgreyL}
2 Kiosk App &
Downloads the latest Kato installer into
\wpath{SelfBooker\textbackslash Documents} and checks it against its
published checksum. Installs it as \wpath{SelfBooker} with the Kiosk ID,
Workstation ID, API key and the \wpath{api.weknow.technology} endpoints,
then checks the version and \wpath{kiosk-identity.json}. If
\wpath{SelfBooker} is signed in, starts Kato for 30 seconds to prove it
runs. &
The download and install, if the latest version is installed with the
same IDs.\\
3 Shell Lock &
What \wpath{SLUtility} does by hand: turns on the Windows Shell Launcher
feature if it is off, makes \wpath{Kiosk.exe} the shell for
\wpath{SelfBooker} (Action~3) and enables Shell Launcher. Then restarts
the machine after 15 seconds. &
Everything, including the restart, if already locked down with Kato
running.\\
\end{tabularx}
\end{center}

\begin{wkwarn}[Read step 1's summary --- a warning does not stop the run]
If the FCC log never shows all three lines, step 1 still ends in
\texttt{SUCCESS}, with a yellow warning and a summary \texttt{Log} line
that says \texttt{checks INCOMPLETE - do not hand the machine on}. The run
carries on and restarts the machine. Fix it with Appendix~\ref{sec:fcc}'s
troubleshooting (\S\ref{sec:fcctrouble}) before you leave site.
\end{wkwarn}

\begin{wkinfo}[Missed the output? It is saved on the machine]
Everything the script prints is also added to
\wpath{C:\textbackslash installation.log}, so it can still be read after
the restart. Reopen the N-central shell in PowerShell mode and run
\wpath{Get-Content C:\textbackslash installation.log -Tail 80}. Each run
starts with \texttt{=== setup\_kiosk on <computer> at <time> ===} and is
added after the earlier ones. The API key is never written to it.
\end{wkinfo}

\begin{wkinfo}[If it stops]
The run ends with \texttt{=== STOPPED at <step> ===}, under that step's
\texttt{FAILED} message and a \texttt{Failed at} line naming what it was
doing. Fix the cause (\S\ref{sec:scripttrouble}), then paste the script
again. Rerunning is safe and quick: steps already done are skipped, so a
rerun on a finished kiosk changes nothing.
\end{wkinfo}

\begin{wkinfo}[The script does not set the printer]
The app installs with the default port, \wpath{COM8}. Step 2's summary
shows the port it ended up with on the \texttt{App} line. If the printer
is elsewhere, correct it in \wpath{kiosk-identity.json} --- see
\S\ref{sec:comport}.
\end{wkinfo}

\paragraph{Running one step on its own.} Each step is also a script of its
own, pasted the same way, with only its own values at the top. Use them to
redo one step, or to build a machine (step~1) before the kiosk record
exists.

\begin{center}
\renewcommand{\arraystretch}{1.3}\small
\begin{tabularx}{\textwidth}{L{40mm} L{34mm} X}
\rowcolor{weknownight}\color{white}\textbf{Script} &
\color{white}\textbf{Values to set} &
\color{white}\textbf{Does}\\
\href{https://raw.githubusercontent.com/vedant-kalbag-wkg/kato-installation-guide/main/scripts/v2/1-FreedomPay.ps1}{\wpath{1-FreedomPay.ps1}} &
\wpath{\$Region} & Step 1 --- Appendix~\ref{sec:fcc}.\\
\rowcolor{weknowgreyL}
\href{https://raw.githubusercontent.com/vedant-kalbag-wkg/kato-installation-guide/main/scripts/v2/2-Kiosk-App.txt}{\wpath{2-Kiosk-App.txt}} &
\wpath{\$KioskId}, \wpath{\$ApiKey} & Step 2 --- the Kato app.\\
\href{https://raw.githubusercontent.com/vedant-kalbag-wkg/kato-installation-guide/main/scripts/v2/3-Shell-Lock.txt}{\wpath{3-Shell-Lock.txt}} &
none & Step 3 --- the lockdown and restart.\\
\end{tabularx}
\end{center}

\begin{wkinfo}[Running it from the N-central repository instead]
Instead of pasting, the script can be stored in N-central once and run
with the values passed in:
\begin{enumerate}
  \item Save the script as \wpath{setup\_kiosk.ps1}. In
        \textbf{Configuration $\rightarrow$ Scheduled Tasks $\rightarrow$
        Script/Software Repository}, select \textbf{Add $\rightarrow$
        Scripting} and upload it. When the script changes, open the item
        and use \textbf{Change} beside the file name.
  \item On the device, open \textbf{Tools $\rightarrow$ Task Execution},
        choose the script, and run it with \textbf{LocalSystem} and these
        command-line parameters:\\
        \wpath{-Region UK -KioskId <id> -ApiKey <key>}
  \item Read the \textbf{Output} pane for
        \texttt{=== ALL 3 STEPS PASSED ===}, or
        \wpath{C:\textbackslash installation.log} after the restart.
\end{enumerate}
N-central keeps a task's command-line parameters, so this stores the API
key in N-central. Regenerate the key if the task record could be seen by
anyone who should not have it.
\end{wkinfo}

\subsection{Option B --- By hand}\label{sec:manual}

\begin{enumerate}
  \item Build the machine with Appendix~\ref{sec:fcc}, signed in as
        \wpath{WKGAdmin}, unless step~1 of the script already passed.
  \item Get the installer onto the machine (\S\ref{sec:getinstaller}).
  \item Install the app, signed in as \wpath{SelfBooker}
        (\S\ref{sec:silent}).
  \item Lock the machine into kiosk mode (\S\ref{sec:lockdown}).
\end{enumerate}

\subsection{Get the installer onto the machine}\label{sec:getinstaller}

\begin{enumerate}
  \item Download the latest \wpath{Kiosk Setup X.Y.Z.exe} from the
        \href{https://github.com/WeKnowTechnology/wkg-electron-app/releases}{app
        releases page}.

  \item Copy it into\\
        \wpath{C:\textbackslash Users\textbackslash SelfBooker\textbackslash Documents}\\
        To copy it remotely, push it from N-central (\S\ref{sec:filepush}).
\end{enumerate}

\subsection{Push a file to kiosks from N-central}\label{sec:filepush}

Use this to add or replace a file on one kiosk or many at once --- a new
installer, an updated config. Two ways; neither needs a remote session.

\paragraph{Option 1 --- File Transfer task (upload the file each time).}
This task is not on a single device's \textbf{Add Task} menu; it is
created at Customer or Service Organization level.
\begin{enumerate}
  \item In N-central, select the \textbf{customer} (not a site or
        device), then \textbf{Configuration $\rightarrow$ Scheduled Tasks
        $\rightarrow$ Add/Delete $\rightarrow$ Add $\rightarrow$ File
        Transfer}. If it is missing, your role needs the
        \textbf{Scheduled Tasks} permission.
  \item Choose \textbf{Use LocalSystem Credentials}, \textbf{Browse} to
        the file (under 500\,MB), and set \textbf{Remote Location} to the
        destination folder, e.g.\
        \wpath{C:\textbackslash Users\textbackslash SelfBooker\textbackslash Documents}.
        Turn on the \textbf{overwrite} option to replace an older copy.
  \item On \textbf{Targets}, tick every kiosk that needs the file.
  \item On \textbf{Schedule}, run it now, then \textbf{Save}. Check each
        kiosk under \textbf{Reports $\rightarrow$ Status $\rightarrow$
        Scheduled Tasks Summary}.
\end{enumerate}
When the file is updated, add a new task with the new file.

\paragraph{Option 2 --- Download script (update the file in one place).}
Host the file at an HTTPS link the kiosks can reach (e.g.\ a SharePoint
or Azure storage download link) and keep that link the same when the
file changes. \wpath{Get-KioskFile.ps1} downloads it onto each kiosk and
replaces any older copy.
\begin{enumerate}
  \item Once: add \wpath{scripts/Get-KioskFile.ps1} from the guide's
        repository to \textbf{Script/Software Repository $\rightarrow$
        Add $\rightarrow$ Scripting}.
  \item At customer level, \textbf{Configuration $\rightarrow$ Scheduled
        Tasks $\rightarrow$ Add/Delete $\rightarrow$ Add $\rightarrow$
        Scripting}: choose the script, \textbf{LocalSystem}, and the
        command-line parameters\\
        \wpath{-Url '<https link>' -Destination '<folder>'}\\
        (add \wpath{-FileName '<name>'} if the link does not end in the
        file name).
  \item Tick the kiosks on \textbf{Targets} and run it now. Each kiosk
        reports \texttt{=== SUCCESS ===} with the file's SHA256, which
        should match on every kiosk.
\end{enumerate}
When the file is updated, replace it at the link and re-run the task.

\subsection{Install the app as \wpath{SelfBooker}}\label{sec:silent}

Done \textbf{signed in as \wpath{SelfBooker}} --- the account the kiosk
runs as.

\begin{enumerate}
  \item Sign in to Windows as \wpath{SelfBooker}.

  \item Open \textbf{PowerShell} in
        \wpath{C:\textbackslash Users\textbackslash SelfBooker\textbackslash Documents}
        and run the silent install, substituting the installer version and
        your values from \S\ref{sec:getkeys}:
\end{enumerate}

\begin{lstlisting}
& ".\Kiosk Setup X.Y.Z.exe" /S `
  /KIOSKID={KIOSK_ID_FROM_COMMAND_CENTRE} `
  /WORKSTATIONID={WKG-POS-XXX} `
  /APIKEY={API_KEY_FROM_COMMAND_CENTRE} `
  /PRINTER=COM8
\end{lstlisting}

\begin{center}
\renewcommand{\arraystretch}{1.3}\small
\begin{tabularx}{\textwidth}{L{34mm} X}
\rowcolor{weknownight}\color{white}\textbf{Switch} &
\color{white}\textbf{Value}\\
\wpath{/S} & Silent install --- no prompts. Always used.\\
\rowcolor{weknowgreyL}
\wpath{/KIOSKID} & The Kiosk ID from the kiosk's Overview tab.\\
\wpath{/WORKSTATIONID} & The machine label, e.g.\ \wpath{WKG-POS-617} ---
the kiosk's serial number.\\
\rowcolor{weknowgreyL}
\wpath{/APIKEY} & The API key you just generated.\\
\wpath{/PRINTER} & The printer's COM port. \wpath{COM8} on a standard
build. If nothing prints, the port is wrong --- see
\S\ref{sec:comport}.\\
\end{tabularx}
\end{center}

\begin{enumerate}\setcounter{enumi}{2}
  \item Confirm the app wrote its identity file and run through the
        checks in \S\ref{sec:identity}. The installer can leave the
        \wpath{dev.} server addresses in it; \S\ref{sec:identity} has the
        command that changes them. If the file is missing, create it
        before continuing --- the app will not connect without it.

  \item Sign out of \wpath{SelfBooker}, then sign in as \wpath{WKGAdmin}
        for the lockdown below.
\end{enumerate}

\subsection{Lock the machine into kiosk mode}\label{sec:lockdown}

The Shell Launcher utility replaces the Windows desktop with the Kato app
for the \wpath{SelfBooker} account, so the machine boots straight into the
kiosk and there is no desktop to escape to. Step~3 of the setup script
does the same thing.

\begin{enumerate}
  \item Signed in as \wpath{WKGAdmin}, run\\
        \wpath{C:\textbackslash Kiosk\textbackslash SSM25\_files\textbackslash Shell-Launcher-Utilities\textbackslash SLUtility.exe}

  \item Select \textbf{Disable User Shell}.

  \item \textbf{Close and reopen} \wpath{SLUtility}. It needs the restart
        to pick up the change --- do not skip this.

  \item Select the \wpath{SelfBooker} user and change the shell path to
        the app, leaving the default \textbf{Action 3}:\\
        \wpath{C:\textbackslash Users\textbackslash SelfBooker\textbackslash AppData\textbackslash Local\textbackslash Programs\textbackslash Kiosk\textbackslash Kiosk.exe}

  \item Select \textbf{Set User Shell}.

  \item Select \textbf{Enable Shell Launcher}.

  \item Sign out of \wpath{WKGAdmin} and sign in as \wpath{SelfBooker}.
        The Kato app should launch full screen with no desktop behind it.
        Reboot once to confirm it comes back the same way.
\end{enumerate}

\begin{wkwarn}[Order matters]
Disable $\rightarrow$ restart the utility $\rightarrow$ set the path
$\rightarrow$ set shell $\rightarrow$ enable. Do not set the shell before
restarting \wpath{SLUtility}.
\end{wkwarn}

\begin{wkinfo}[Done when]
The machine reboots straight into the Kato app as \wpath{SelfBooker}, with
no visible desktop, taskbar or file explorer.
\end{wkinfo}

%=====================================================================
% 4  FINAL CHECKLIST
%=====================================================================
\wksection{Final checklist}\label{sec:checklist}

Run through this before leaving site.

\begin{center}
\renewcommand{\arraystretch}{1.45}\small
\begin{tabularx}{\textwidth}{L{8mm} L{14mm} X}
\rowcolor{weknownight}\color{white} & \color{white}\textbf{Who} &
\color{white}\textbf{Check}\\
\emptybox & AM & Kiosk record has the correct name, serial and hardware
model.\\
\rowcolor{weknowgreyL}
\emptybox & AM & All prep checks (hardware, firmware, software) cleared.\\
\emptybox & AM & Installation record created and linked to the location;
kiosk no longer \emph{In Stock}.\\
\rowcolor{weknowgreyL}
\emptybox & AM & Kiosk link sent to IT.\\
\emptybox & IT & Setup script ended in \texttt{ALL 3 STEPS PASSED}, and
step~1's \texttt{Log} line is not \texttt{INCOMPLETE} --- or, by hand,
Appendix~\ref{sec:fcc}'s checks all passed.\\
\rowcolor{weknowgreyL}
\emptybox & IT & \wpath{kiosk-identity.json} present and correct.\\
\emptybox & IT & API key not saved anywhere --- Notepad closed without
saving, nothing in the ticket or chat.\\
\rowcolor{weknowgreyL}
\emptybox & IT & Machine reboots straight into the Kato app --- no
desktop.\\
\emptybox & IT & Printer prints, scanner scans, sensors respond.\\
\rowcolor{weknowgreyL}
\emptybox & IT & Test payment taken on the PDQ and refunded.\\
\end{tabularx}
\end{center}

%=====================================================================
% 5  TROUBLESHOOTING
%=====================================================================
\wksection{When things go wrong}\label{sec:trouble}

\begin{wkinfo}[The command boxes in this guide]
Unless a box says otherwise, its commands are for the N-central remote
connection's shell in PowerShell mode (\S\ref{sec:ncentral}). Copy the
whole box and paste it. Each line is a complete command, so a box also
works pasted one line at a time, and nothing in it waits for input.
\end{wkinfo}

\begin{center}
\renewcommand{\arraystretch}{1.35}\small
\begin{tabularx}{\textwidth}{L{50mm} X}
\rowcolor{weknownight}\color{white}\textbf{Symptom} &
\color{white}\textbf{First thing to try}\\
Machine boots to the Windows desktop, not the app &
Paste step~3 again (\S\ref{sec:ncentral}). By hand, the \wpath{SLUtility}
sequence was not followed: redo it in order, restarting the utility
between \emph{Disable User Shell} and setting the path.\\
\rowcolor{weknowgreyL}
You need the Windows desktop back, e.g.\ for maintenance &
Turn Shell Launcher off (\S\ref{sec:unlock}).\\
App launches but will not connect &
Work through the \wpath{kiosk-identity.json} checks in
\S\ref{sec:identity} --- dev server URLs, a stray backtick on a value, a
wrong Kiosk ID, or an API key regenerated after the install.\\
\rowcolor{weknowgreyL}
Setup script ends in \texttt{STOPPED}, or step~1 warns &
The message names the cause --- see \S\ref{sec:scripttrouble}.\\
The machine restarted before you read the script's output &
Run \wpath{Get-Content C:\textbackslash installation.log -Tail 80}.\\
\rowcolor{weknowgreyL}
Installer runs but nothing appears &
By hand: the switches are case-sensitive and the backtick line
continuations must be the last character on each line. Retype the command
rather than pasting it from a chat client, or use the setup script
(\S\ref{sec:ncentral}).\\
Printer does nothing &
The configured COM port is wrong. Find the port the printer is actually
on and correct it in \wpath{kiosk-identity.json} --- no reinstall needed
(\S\ref{sec:comport}).\\
\rowcolor{weknowgreyL}
Prep checks will not clear &
They are marked manually on the kiosk \textbf{Edit} tab --- they do not
clear themselves when the underlying condition is met.\\
Payment fails on the test transaction &
That is a machine-build issue, not a configuration one. See
Appendix~\ref{sec:fcc}, and its troubleshooting table
(\S\ref{sec:fcctrouble}).\\
\end{tabularx}
\end{center}

\subsection{When the setup script stops or warns}\label{sec:scripttrouble}

The line under \texttt{FAILED} says what went wrong, and the summary's
\texttt{Failed at} line says what the step was doing. Fix it, then paste
the setup script again --- steps already done are skipped. To restart the
machine from the N-central shell, run \wpath{Restart-Computer -Force}.

\medskip\noindent Values the script rejects:

\begin{center}
\renewcommand{\arraystretch}{1.35}\small
\begin{tabularx}{\textwidth}{L{62mm} X}
\rowcolor{weknownight}\color{white}\textbf{Message} &
\color{white}\textbf{Fix}\\
\texttt{Set ... at the top of the script before running} &
A value is still a placeholder. Set it in Notepad and paste again.\\
\rowcolor{weknowgreyL}
\texttt{Region must be one of ...} &
Use \wpath{UK}, \wpath{DE}, \wpath{ES} or \wpath{PT}, in capitals.\\
\texttt{KioskId is malformed}, \texttt{ApiKey is malformed} &
Mistyped or partly copied. Copy it again from Command Centre --- an API
key starts \wpath{kcs\_}. Generate a new key if the old one is gone.\\
\rowcolor{weknowgreyL}
\texttt{WorkstationId is malformed} &
The computer name is not a workstation label. Use Option~B
(\S\ref{sec:manual}).\\
\end{tabularx}
\end{center}

\wkneed{10\baselineskip}\noindent Problems during the run:

\begin{center}
\renewcommand{\arraystretch}{1.35}\small
\begin{tabularx}{\textwidth}{L{62mm} X}
\rowcolor{weknownight}\color{white}\textbf{Message} &
\color{white}\textbf{Fix}\\
\texttt{FCC installer still running after 10 min} &
FreedomPay's installer is waiting on a window nobody can see. Install FCC
by hand (Appendix~\ref{sec:fcc}), then paste again --- step~1 skips the
install.\\
\rowcolor{weknowgreyL}
\texttt{FCC install failed} &
The message names the install log. Raise it with the WKG Technology team,
with that log.\\
\texttt{... is not signed in and does not sign in automatically},
\texttt{No stored sign-in password ...} &
Sign in to Windows as \wpath{SelfBooker} on the machine, then paste
again.\\
\rowcolor{weknowgreyL}
\texttt{Downloaded installer does not match ...} &
The download was damaged. Paste again.\\
\texttt{Config mismatch: ...} &
The identity file does not hold the values given. Check it
(\S\ref{sec:identity}), then paste again --- step~2 reinstalls when the
IDs differ.\\
\rowcolor{weknowgreyL}
\texttt{Kato closed within 30s of starting} &
Work through the identity file checks (\S\ref{sec:identity}).\\
\texttt{Shell Launcher was just turned on and needs a restart} &
Restart the machine, then paste again. Steps~1 and~2 are skipped.\\
\rowcolor{weknowgreyL}
Step~1 warns \texttt{FCC log has no Lane Open ...}, and its summary says
\texttt{checks INCOMPLETE} &
The card terminal did not confirm the build. See
\S\ref{sec:fcctrouble} before you leave site.\\
\end{tabularx}
\end{center}

\subsection{Getting the Windows desktop back}\label{sec:unlock}

This turns Shell Launcher off, so \wpath{SelfBooker} gets the normal
Windows desktop at the next sign-in:

\begin{lstlisting}
$sl = @{ Namespace = 'root\standardcimv2\embedded'; ClassName = 'WESL_UserSetting' }
Invoke-CimMethod @sl -MethodName SetEnabled -Arguments @{ Enabled = $false }
\end{lstlisting}

Then restart the machine. To lock it again, paste step~3
(\wpath{3-Shell-Lock.txt}) or the whole setup script again; by hand, use
\wpath{SLUtility}'s \textbf{Enable Shell Launcher}.

\subsection{Checking \wpath{kiosk-identity.json}}\label{sec:identity}

Most \emph{will not connect} faults are in this file (path in
\S\ref{sec:paths}). To see it with the API key hidden, run:

\begin{lstlisting}
$f = 'C:\Users\SelfBooker\AppData\Roaming\WKG Kiosk\kiosk-identity.json'
$j = Get-Content $f -Raw | ConvertFrom-Json
if ($j.apiKey) { $j.apiKey = $j.apiKey.Substring(0,4) + '****' }
$j | Format-List
\end{lstlisting}

A correct file looks like this:

\begin{lstlisting}
{
  "kioskId": "{KIOSK_ID_FROM_COMMAND_CENTRE}",
  "apiBaseUrl": "https://api.weknow.technology/api",
  "apiKey": "{API_KEY_FROM_COMMAND_CENTRE}",
  "workstationId": "WKG-POS-XXX",
  "printerDestination": "COM8",
  "fleetSocketUrl": "wss://api.weknow.technology"
}
\end{lstlisting}

\begin{center}
\renewcommand{\arraystretch}{1.35}\small
\begin{tabularx}{\textwidth}{L{40mm} X}
\rowcolor{weknownight}\color{white}\textbf{Field} &
\color{white}\textbf{Check}\\
\wpath{apiBaseUrl} &
Must be \wpath{https://api.weknow.technology/api} --- \textbf{not}
\wpath{https://dev.weknow.technology/api}.\\
\rowcolor{weknowgreyL}
\wpath{fleetSocketUrl} &
Must be \wpath{wss://api.weknow.technology} --- \textbf{not}
\wpath{wss://dev.weknow.technology}.\\
\wpath{kioskId} &
Matches the Kiosk ID on the kiosk's Overview tab, and does \textbf{not end
in a backtick} (\wpath{`}).\\
\rowcolor{weknowgreyL}
\wpath{apiKey} &
Matches the key generated for this kiosk, and does \textbf{not end in a
backtick}. Only its first four characters are shown.\\
\wpath{workstationId} &
The machine label, e.g.\ \wpath{WKG-POS-617} --- not the Kiosk ID.\\
\rowcolor{weknowgreyL}
\wpath{printerDestination} &
The printer's COM port (\S\ref{sec:comport}).\\
\end{tabularx}
\end{center}

\begin{wkwarn}[Why a backtick ends up in a value]
In the manual install command, each line ends in a backtick to continue
it. If a space follows the backtick, or the command is pasted onto one
line, PowerShell passes the backtick through as part of the value. Delete
it from the file.
\end{wkwarn}

\paragraph{Fixing it.} There is no need to reinstall. To change the
\wpath{dev.} server addresses to the live ones, run the box below. It
keeps a copy of the file as \wpath{kiosk-identity.json.bak}, then shows the
result with the API key hidden.

\begin{lstlisting}
$f = 'C:\Users\SelfBooker\AppData\Roaming\WKG Kiosk\kiosk-identity.json'
Copy-Item $f "$f.bak" -Force
$t = (Get-Content $f -Raw) -replace '://dev\.weknow\.technology', '://api.weknow.technology'
[IO.File]::WriteAllText($f, $t, (New-Object Text.UTF8Encoding $false))
$j = Get-Content $f -Raw | ConvertFrom-Json; $j.apiKey = 'kcs_****'; $j | Format-List
\end{lstlisting}

The copy holds the API key too, so delete it once Kato connects:

\begin{lstlisting}
Remove-Item 'C:\Users\SelfBooker\AppData\Roaming\WKG Kiosk\kiosk-identity.json.bak'
\end{lstlisting}

Fix any other wrong value in the file in Notepad and save it, or paste
step~2 (\wpath{2-Kiosk-App.txt}) again with the right values --- it
reinstalls when the Kiosk ID, Workstation ID or API key differ. If the API
key no longer works at all, generate a new one in Command Centre and use
that. Then restart the machine so Kato reads the file again.

\subsection{Finding the printer's COM port}\label{sec:comport}

\begin{wkwarn}[Device Manager is not the answer here]
The app writes \textbf{straight to the COM port} --- there is no printer
driver in between. The port the printer actually answers on does
\emph{not} always match what Device Manager shows against it, so read the
ports, then test them.
\end{wkwarn}

\paragraph{1. List the ports Windows can see.} In PowerShell:

\begin{lstlisting}
# Just the port names
[System.IO.Ports.SerialPort]::GetPortNames()

# ...with the device name against each, which narrows the field
$ports = Get-CimInstance Win32_PnPEntity | Where-Object { $_.Name -match 'COM\d+' }
$ports | Select-Object Name, DeviceID
\end{lstlisting}

Or, from Command Prompt (the N-central shell's \textbf{CMD} mode),
\wpath{mode} lists every COM device with its current settings:

\begin{lstlisting}
mode
\end{lstlisting}

\paragraph{2. Test each port until the printer reacts.} Writing anything
to the right port makes the printer spit out a few lines of nonsense ---
that is the confirmation. Run this from \textbf{Command Prompt}, one port
at a time, working through the list from step 1:

\begin{lstlisting}
DIR > COM3
DIR > COM4
DIR > COM8
\end{lstlisting}

From PowerShell, hand the redirect to \wpath{cmd} instead, so it reaches
the device rather than creating a file:

\begin{lstlisting}
cmd /c "DIR > COM8"
\end{lstlisting}

How to read the result:

\begin{center}
\renewcommand{\arraystretch}{1.3}\small
\begin{tabularx}{\textwidth}{L{54mm} X}
\rowcolor{weknownight}\color{white}\textbf{What happens} &
\color{white}\textbf{Means}\\
Printer prints garbage &
That is the port. Note it and move to step 3.\\
\rowcolor{weknowgreyL}
Command returns silently, nothing prints &
The port exists but the printer is not on it. Try the next one.\\
\emph{Access is denied} &
Something already holds the port --- usually the kiosk app. Close it and
try again.\\
\rowcolor{weknowgreyL}
\emph{The system cannot find the path specified} &
No such port. Try the next one.\\
\end{tabularx}
\end{center}

\paragraph{3. Set the port in \wpath{kiosk-identity.json}.} There is
\textbf{no need to rerun the silent install}. Open
\wpath{kiosk-identity.json}, set the printer port to the one you
confirmed, save, then restart the app and print a test receipt.

\begin{wkinfo}[Note it on the installation record]
If the printer turned out to be on a non-standard port, put it in the
\textbf{Notes} field of the installation record. It saves the next
engineer repeating this.
\end{wkinfo}

\clearpage
\appendix
%=====================================================================
% APPENDIX A -- BUILDING THE MACHINE (FreedomPay FCC)
%=====================================================================
\wkrolesection{Building the machine --- FreedomPay FCC}{IT}\label{sec:fcc}

Installing FreedomPay Freeway Commerce Connect (FCC), applying the WeKnow
configuration and PAL bundle, and confirming the card terminal works. Done
\textbf{once per machine}, before the app is installed.

\begin{wkinfo}[The setup script does all of this for you]
Step~1 of the Part~2 setup script (\S\ref{sec:ncentral}) does this whole
appendix and its checks, for the region you set. \wpath{1-FreedomPay.ps1}
runs that step on its own. Use this appendix to build a machine by hand,
or to check or fix what the script did.
\end{wkinfo}

\subsection{What you need}

Do everything in this appendix signed in as \wpath{WKGAdmin}.

\begin{center}
\renewcommand{\arraystretch}{1.35}\small
\begin{tabularx}{\textwidth}{L{44mm} X}
\rowcolor{weknownight}\color{white}\textbf{Item} &
\color{white}\textbf{Where it comes from}\\
FCC installer &
Internal mirror: \wpath{https://dev.weknowlondon.uk/SSM/freedompay/}.
The scripts download it from
\wpath{https://wkgdevstore.blob.core.windows.net/kiosk/FCC/}.\\
\rowcolor{weknowgreyL}
Config + PAL bundle for the region &
UK from the same mirror. The scripts take UK, DE, ES or PT from the same
place as the installer.\\
PDQ (card terminal) &
Connected to the machine --- you cannot verify the build without it.\\
\end{tabularx}
\end{center}

\subsection{Install FreedomPay FCC}

Skip to \S\ref{sec:fccverify} if FCC is already installed on this machine and
you only need to confirm it is healthy.

\begin{enumerate}
  \item Run \wpath{FreewayCommerceConnect.exe} as an admin user.

  \item \textbf{Skip the activation step} of the Remote Update Service
        (RUS).

  \item Choose a \textbf{Custom} installation, ticking the features shown
        below.
\end{enumerate}

\screenshot[0.50]{fcc-custom-features}{Custom Setup feature tree --- \textbf{Freeway POI Drivers} and \textbf{Freeway Server Service} ticked, everything else left clear.}

\begin{center}
\renewcommand{\arraystretch}{1.3}\small
\begin{tabularx}{\textwidth}{L{30mm} X}
\rowcolor{weknownight}\color{white}\textbf{Purpose} &
\color{white}\textbf{Path}\\
Binaries \& config &
\wpath{C:\textbackslash Program Files (x86)\textbackslash FreedomPay\textbackslash FreewayCommerceConnect}\\
\rowcolor{weknowgreyL}
Logs &
\wpath{C:\textbackslash ProgramData\textbackslash FreedomPay\textbackslash Freeway Commerce Connect\textbackslash log}\\
\end{tabularx}
\end{center}

\subsection{Apply the WeKnow configuration}

\begin{enumerate}
  \item Open the Registry Editor and \textbf{delete} the key\\
        \wpath{HKEY\_LOCAL\_MACHINE\textbackslash SOFTWARE\textbackslash Freedompay\textbackslash Atlas\textbackslash PAL}

  \item Copy the UK \wpath{config} and \wpath{PAL} folders into\\
        \wpath{C:\textbackslash Program Files (x86)\textbackslash FreedomPay\textbackslash FreewayCommerceConnect}

  \item Restart the services in this order:
        \wpath{FreewayServerService} $\rightarrow$ \wpath{FreewayClientService}
\end{enumerate}

\subsection{Verify the build}\label{sec:fccverify}

Do not hand the machine on until all four checks below pass.

\begin{enumerate}
  \item In the Windows Services app, confirm both
        \wpath{FreewayClientService} and \wpath{FreewayServerService}
        show \textbf{Running}.

  \item Open the newest log file in\\
        \wpath{C:\textbackslash ProgramData\textbackslash FreedomPay\textbackslash Freeway Commerce Connect\textbackslash log}\\
        and confirm the latest timestamp for \textbf{Lane Open}.

  \item In the same log, confirm the regional configuration package
        reached the terminal --- look for \emph{Successfully pushed}.
\end{enumerate}

\screenshot[0.94]{fcc-log-palpush}{\emph{Successfully pushed} --- the regional config package reached the terminal.}

\begin{enumerate}\setcounter{enumi}{3}
  \item Confirm the terminal accepted and wrote it --- look for a
        \textbf{PalWriteSucceeded} event.
\end{enumerate}

\screenshot[0.94]{fcc-log-palwrite}{\emph{PalWriteSucceeded} --- the terminal accepted and wrote the package.}

\begin{wkinfo}[Done when]
Both services are running, and the log shows \textbf{Lane Open}, a
successful package push and a \textbf{PalWriteSucceeded} at the latest
timestamp. The machine is now ready for Part~2 (\S\ref{sec:part2}).
\end{wkinfo}

\subsection{When the build goes wrong}\label{sec:fcctrouble}

\begin{center}
\renewcommand{\arraystretch}{1.35}\small
\begin{tabularx}{\textwidth}{L{50mm} X}
\rowcolor{weknownight}\color{white}\textbf{Symptom} &
\color{white}\textbf{First thing to try}\\
No \emph{Lane Open} in the log &
The terminal is not being detected. Check the cable, then restart the
services with the first box below --- it waits up to a minute for a new
\emph{Lane Open}. Still nothing: power-cycle the terminal and run it
again.\\
\rowcolor{weknowgreyL}
No \emph{Lane Open} after the setup script built the machine &
Check the PDQ driver is installed. The script's silent FCC install copies
the driver's installer, \wpath{FreewayMsrDriverExe.exe}, into the FCC
folder (\wpath{C:\textbackslash Program Files (x86)\textbackslash FreedomPay\textbackslash FreewayCommerceConnect})
but does not run it. Signed in as \wpath{WKGAdmin} with the PDQ plugged
in, run it from there. Then clear the PAL record (second box below) and
paste step~1 again, so the package is pushed to the terminal.\\
No \emph{Successfully pushed} or \emph{PalWriteSucceeded} &
Confirm the \wpath{Atlas\textbackslash PAL} registry key was deleted
\emph{before} the config and PAL folders were copied in, then restart the
services in order (\wpath{FreewayServerService} $\rightarrow$ \wpath{FreewayClientService}).
After the setup script: clear the PAL record and paste step~1 again.\\
\rowcolor{weknowgreyL}
Step~1 says the PAL package is \texttt{already taken}, but the terminal
does not have it &
Step~1 records the last package it ran and does not run it again. Clear
the record (second box below), then paste step~1 again.\\
Services will not stay running &
A partially copied bundle or an unsupported version of FCC is the usual cause. Reinstall a known good version of FCC. Recopy the \wpath{config}
and \wpath{PAL} folders and restart the services in order.\\
\end{tabularx}
\end{center}

\paragraph{Restart the FCC services and wait for \emph{Lane Open}.}
This restarts the server service, then the client service --- their
Windows names are \wpath{FCCServerSvc} and \wpath{FCCClientSvc}. It shows
both are running, then watches the FCC log for up to 60 seconds:

\begin{lstlisting}
$log = 'C:\ProgramData\FreedomPay\Freeway Commerce Connect\log'
$n = { @(Get-ChildItem $log -File | Select-String -Pattern 'Lane Open').Count }
$before = & $n
Restart-Service FCCServerSvc -Force; Restart-Service FCCClientSvc -Force
Get-Service FCCServerSvc, FCCClientSvc | Format-Table Name, Status
$ok = $false
for ($i = 0; $i -lt 30 -and -not $ok; $i++) { Start-Sleep 2; $ok = (& $n) -gt $before }
if ($ok) { 'Lane Open - terminal detected' } else { 'No new Lane Open after 60s' }
\end{lstlisting}

\emph{No new Lane Open after 60s}: check the PDQ cable and driver, then
power-cycle the terminal.

\paragraph{Clear step~1's PAL record.} The next run of step~1 then runs
the PAL package again:

\begin{lstlisting}
reg delete "HKLM\SOFTWARE\WeKnow\Kiosk" /v PalPackage /f /reg:64
\end{lstlisting}

\begin{wkinfo}[Still stuck?]
Zip the whole \wpath{log} folder, note the workstation serial and the
terminal serial, and raise it with the WKG Technology team before
changing anything else.
\end{wkinfo}


%=====================================================================
% APPENDIX B -- ONE-PAGE QUICK REFERENCE
%=====================================================================
\clearpage
\wksection{Quick reference}\label{sec:quickref}

One page for someone who has done this before. Work down your column;
the full detail is in the sections above.

\begin{center}
\begin{minipage}[t]{0.475\textwidth}
\colorbox{weknowsky}{\parbox{\dimexpr\textwidth-2\fboxsep}{%
  \vspace{2pt}\centering\bfseries\color{weknownight}
  Part 1 --- Account Manager\vspace{2pt}}}
\vspace{4pt}
\footnotesize
\begin{enumerate}[leftmargin=6mm,itemsep=3pt]
  \item \emptybox\ Collect the location information (\S1.1). Skip RPS
        customer code and outlet code --- legacy.
  \item \emptybox\ Create the location; copy the \textbf{Location ID}.
  \item \emptybox\ Create groups if needed. Kiosk groups can span sites.
  \item \emptybox\ \textbf{Add Kiosk} --- name, serial, hardware model,
        firmware.
  \item \emptybox\ Assign kiosk groups on the \textbf{Edit} tab.
        Bottom of the list wins a conflict --- most specific last.
  \item \emptybox\ Set the \textbf{Store ID} and \textbf{Terminal ID}.
  \item \emptybox\ Clear the prep checks --- only after IT confirm
        hardware and FreedomPay.
  \item \emptybox\ Create the \textbf{installation record}.
  \item \emptybox\ Send IT the kiosk's Command Centre \textbf{link} ---
        not sensitive. Mention the COM port if not \wpath{COM8}.
\end{enumerate}
\end{minipage}\hfill
\begin{minipage}[t]{0.475\textwidth}
\colorbox{weknowgreyL}{\parbox{\dimexpr\textwidth-2\fboxsep}{%
  \vspace{2pt}\centering\bfseries\color{weknownight}
  Part 2 --- IT\vspace{2pt}}}
\vspace{4pt}
\footnotesize
\begin{enumerate}[leftmargin=6mm,itemsep=3pt]
  \item \emptybox\ PDQ plugged in. Computer name is the Workstation ID.
        \wpath{SelfBooker} signed in or signs in automatically.
  \item \emptybox\ From the kiosk link: copy the \textbf{Kiosk ID};
        generate the \textbf{API key}. Use it at once --- never save it.
  \item \emptybox\ N-central remote connection: \textbf{System Shell},
        PowerShell mode.
  \item \emptybox\ Open \wpath{setup\_kiosk.ps1}; copy it into Notepad;
        set \wpath{\$Region}, \wpath{\$KioskId}, \wpath{\$ApiKey}.
  \item \emptybox\ Copy it all; close Notepad without saving; paste it
        into the shell.
  \item \emptybox\ Wait for \texttt{ALL 3 STEPS PASSED}. Step~1's
        \texttt{Log} line must not say \texttt{INCOMPLETE}.
  \item \emptybox\ It restarts into Kato. Missed the output? It is in
        \wpath{C:\textbackslash installation.log}.
  \item \emptybox\ \texttt{STOPPED}? Fix it (\S\ref{sec:scripttrouble})
        and paste again --- done steps are skipped.
  \item \emptybox\ By hand instead: Appendix~\ref{sec:fcc}, the silent
        install below, then \wpath{SLUtility} (\S\ref{sec:lockdown}).
  \item \emptybox\ Test print, test scan, test payment.
\end{enumerate}
\end{minipage}
\end{center}

\vspace{6pt}

\begin{center}
\renewcommand{\arraystretch}{1.2}\footnotesize
\begin{tabularx}{\textwidth}{L{34mm} X}
\rowcolor{weknownight}\color{white}\textbf{Value} &
\color{white}\textbf{Where IT get it}\\
Kiosk link & From the AM, in the ticket or chat. Not sensitive.\\
\rowcolor{weknowgreyL}
Kiosk ID & UUID on the kiosk's Overview tab.\\
API key & Generated by IT on the API keys tab at install time. Shown once;
paste straight in, never saved or sent.\\
\rowcolor{weknowgreyL}
Workstation ID & The kiosk's serial number, \wpath{WKG-POS-XXX}.\\
Printer COM port & \wpath{COM8} unless the AM says otherwise.\\
\end{tabularx}
\end{center}

\vspace{4pt}
\noindent{\footnotesize\textbf{Silent install by hand} --- run in
PowerShell, signed in as \wpath{SelfBooker}, from the folder holding the
installer (\S\ref{sec:silent}). Then check \wpath{kiosk-identity.json}
--- IDs, API key, no trailing backtick, \wpath{api.} not \wpath{dev.}
URLs (\S\ref{sec:identity}).}

\begin{lstlisting}
& ".\Kiosk Setup X.Y.Z.exe" /S /KIOSKID={KIOSK_ID} /WORKSTATIONID={WKG-POS-XXX} /APIKEY={API_KEY} /PRINTER=COM8
\end{lstlisting}

\begin{center}
\renewcommand{\arraystretch}{1.2}\footnotesize
\begin{tabularx}{\textwidth}{L{34mm} X}
\rowcolor{weknownight}\color{white}\textbf{Path} &
\color{white}\textbf{Where}\\
Kato app &
\wpath{C:\textbackslash Users\textbackslash SelfBooker\textbackslash AppData\textbackslash Local\textbackslash Programs\textbackslash Kiosk\textbackslash Kiosk.exe}\\
\rowcolor{weknowgreyL}
Shell Launcher utility &
\wpath{C:\textbackslash Kiosk\textbackslash SSM25\_files\textbackslash Shell-Launcher-Utilities\textbackslash SLUtility.exe}\\
Find the printer port &
List with \wpath{mode}, then \wpath{DIR > COM3} from Command Prompt.
Fix in \wpath{kiosk-identity.json} --- no reinstall.\\
\rowcolor{weknowgreyL}
Setup script log &
\wpath{C:\textbackslash installation.log} --- read it with
\wpath{Get-Content C:\textbackslash installation.log -Tail 80}.\\
\end{tabularx}
\end{center}

\end{document}
